\documentclass{article}
\usepackage{hyperref}
\usepackage{float}
\usepackage{csquotes}
\usepackage[style=iso]{datetime2}
\usepackage[usenames,dvipsnames]{color}
\usepackage{booktabs}
  \setlength\heavyrulewidth{0.20ex}
  \setlength\cmidrulewidth{0.10ex}
  \setlength\lightrulewidth{0.10ex}

\usepackage[font=normalsize,labelfont={bf}]{caption}
  \captionsetup[table]{aboveskip=3pt}

\hypersetup{
  colorlinks=true,
  linkcolor=blue,
  filecolor=magenta,
  urlcolor=magenta,
}
% TO-DO
% add in an a5 pdf version? (undecided if this is crucial? i mean i DO want to make the pdfs prettier....)
% add auto loc func for available loc when setting up loc

\usepackage{graphicx}
\graphicspath{ {../export/img/} }
\renewcommand{\abstractname}{免责声明}
\title{女性化HRT实用指南}
\author{\href{https://katea.gay/}{Katie Tightpussy}}
\date{\today}
\setcounter{section}{-1}
\setcounter{figure}{-1}
\urlstyle{same}

\begin{document}


\maketitle
\tableofcontents
\begin{abstract}
  我不是医生。我也不在医学领域工作。我在任何意义上都不是医疗专业人员。我只是一个普通人，基于我自身所受的教育程度与个人经验，提供普通人的观点。以下所有信息和论断，都应当被视为意见，而不是事实陈述或医疗建议。本指南在科学研究力有不逮之处，优先参考社群内部长期实践中形成的。简单来说，别冲我发火。
\end{abstract}

\section*{语言}
\addcontentsline{toc}{section}{语言}

本文档目前提供以下语言版本：

LANGUAGE-CODE-DOT-PGHRT-DOT-DIY

\noindent\textbf{如果你有兴趣参与翻译或制作其他版本，欢迎联系我！}



\section{前言}

这是一份活文档。

我写它的目的是系统性地整理我自己关于女性化激素替代治疗（HRT）的看法和观点，因为我认为现有的各类社群和Wiki并不是很好用。它们当然是很宝贵的资源，但我认为这些 Wiki 对那些想要清晰而可执行的指南的人来说，实用性很差。它们更适合那些想把所有有关的生物学过程和所有图表都学一遍的人。而我想做的，是一份详尽但快速查阅的参考指南。用尽量简化、直接的方式，回答我这些年来最常被问到的问题：关于如何安全、有效地进行 HRT。我的目标，是让消除这种拯救生命的医疗手段所笼罩的迷思，既面向正在考虑开始 HRT 的MtF，也面向已经在进行中的MtF。因此，我默认你已经对HRT的效果有了基本的了解。如果你还不太熟悉：HRT的作用真的非常多，比你想象的还要多得多。它真的很棒。\textbf{改变你的性别是一件非常酷，也非常有趣的事情。强烈推荐。}你值得拥有高质量的过渡医疗服务（transition healthcare），你也有能力为自己做出最合适的决定。我希望这份文档能成为你决策过程中的一个有用工具，如果你想继续深入学习，它也可以作为起点。

还有一条忠告。一定不要混跨性别相关的subreddits。相信我。至少远离 r/mtf，这个社区烂透了。这地方既不健康，也不能提供什么靠谱的经验。你会在之后的很多年里，一点点把那些烂到脑子里的坏观念往外抠。这是我能给你的最好的建议。

至于男同胞们，这份指南里面有些内容和你们依然息息相关，虽然目标和结果，很明显，存在显著差异。\href{https://docs.google.com/document/d/1DXFxzN0XTudPZez\_SO61fpqncRLPH\_Be\_QG\_8Pcz9LU/edit?tab=t.0}{这份男性化HRT指南}\textcolor{red}{(警告：Google Docs链接)}看起来不错，不过我没有一行一行检查过，所以请你们自行判断。总之，他们也该写一份 tboy 版的 Katie Tightpussy。比如叫 Oliver·格调什么之类的。或者Xavier。

\textbf{如果你愿意捐款支持这个项目，} \href{https://cash.app/Katitties}{CashApp}、\href{https://ko-fi.com/katitties}{Ko-Fi}、\href{https://account.venmo.com/u/katitties}{Venmo}都可以！非常感谢！

\subsection*{如何使用本文档}

本文档按线性结构编排，以一连串问答的形式展开。大体上每一个问题和章节都会自然过渡到下一个。我建议你从头到尾按顺序读，因为这样往往能顺带回答一些你自己还没意识到的问题。不得不说，这篇指南确实很长。如果你想分多次慢慢读也是完全没问题的。

你也可以通过目录按需跳转到某个具体的章节或问题，尤其是适合读完查阅时使用。我推荐你把本页面/文档保存下来，这样当你对自己的HRT有任何疑问时，都可以翻出来查阅。一次性吸收这么多内容不容易，所以就算现在看不太懂也完全没关系！

\noindent\textbf{\href{PDF_LINK}{本文档亦提供PDF形式。请务必下载一份。}}

\noindent\href{https://pghrt.diy/pghrtgretchensversion.txt}{或者，如果你更偏好那种上世纪末风格的纯文本体验，也可以阅读这个 .txt 版本。}不过它不会保持最新版就是了。


\section*{献词}
\addcontentsline{toc}{section}{献词}

本文献给我们所有没能活下来的姐妹。愿我们把她们的火炬继续带到明天。

\section{引言}

\subsection{服用雌激素安全吗？}

在现代生物同源激素（bioidentical hormones）条件下，HRT 几乎已经不可能再更安全了。你所做的事情，本质上只是把身体所依赖的主「燃料」换了一种，并重新调整你体内本来就存在的激素平衡。即便在优化细节变得复杂的地方，改变生理结构这件事本身的容错率也非常高。身体是可塑的。你也一定能调整到最适合你自己的状态。

\subsection{我应该选择哪种雌激素给药途径？}

注射。总体而言，注射是最不容易出错、最简单、最稳定、最安全、也最便宜的 HRT 形式。对一部分人来说，注射甚至会变成一种值得期待的仪式。对另一部分人来说，它会变得相当有趣。

\noindent\underline{\textbf{但请记住一句话：不管是怎样服用雌激素，都比不服用雌激素要好。}}

\subsection{为什么你不推荐口服、贴片或凝胶？}

最主要的原因是：这三种方式都有注射所没有的重大缺陷。它们不是没用，而是你不该被迫去忍受这些明显的劣势。我再强调一遍：\textbf{所有 HRT 形式都可能产生令人满意的效果}，但这不代表所有 HRT 形式都是等价的，更不代表它们都是好的选择。

\subsection{不同给药途径或形式下，雌激素剂量是等效的吗？}

不是。这一点重要到，我甚至没有把它丢进第\ref{MM}章《迷思与杂项》里。\textbf{雌激素的剂量，是绝对不能在不同类型或形式之间直接比较的。}1 mg 的这个，绝对不等于 1 mg 的那个。不同类型和形式的雌激素，在其身体吸收雌激素的速率——也就是所谓的生物利用度（bioavailability）、吸收速度、以及半衰期上都存在差异。

\subsection{什么是半衰期？}

简单来说，半衰期指的就是一种物质减半所需的时间。当在HRT的语境下，半衰期决定了一次给药在你体内能维持多久，以及你需要多频繁地用药。这就是你的激素周期。这个周期往往是一条曲线：激素水平上升，达到峰值，然后下降。理解雌激素水平与时间关系的曲线是非常重要的。

\subsection{口服药物有什么问题？}

最主要的一个问题是，口服药物它们会增加长期血栓和肝脏凝血风险。和吞服相比，舌下含服或颊黏膜含服（「含糖」一词的来源——译者注）可以在一定程度上绕开肝脏首过代谢，从而降低一定的风险。但你仍然会不可避免地会咽下一部分药物。因此可以合理地认为，上述风险在一定程度上仍然存在。虽然这么说，绝对风险还是很低的（对乙酰氨基酚/扑热息痛对肝脏的风险甚至都比雌激素的要高一整个数量级）。但是，尼古丁是会进一步放大雌激素相关风险的。详细请参见问题 \ref{11-2}。

除此以外，口服药物还有一些其他的问题，主要有二。其一，它们半衰期很短，生物利用度也不是很高。其二，服用雌激素大多时候需要配套抗雄激素。前者使口服在单一用药方案（monotherapy）方面几乎比不上注射。而后者则往往根据不同的抗雄激素带来一大堆潜在副作用（参见第\ref{AA}章《抗雄激素》）。这两者共同增加了方案中的不确定性，使得糟糕的用药方案及其副作用（如精力和性欲下降、或者见效更慢）也更普遍。此外，口服药物也很难囤积。在一些市场中，口服药物的价格甚至比注射用药物价格还要高。还请注意，从国外渠道大量进口口服药物可能会触犯海关规定，具体后果取决于你所在国家的法律，可能包括货物被扣押、经济损失，甚至潜在的法律风险。\textbf{如果有人问起来，你甚至都不知道这些药物是谁订购的。}

不论你出于什么原因使用口服药物，请使用 4-8 mg，舌下含服，并将剂量分散到一整天内。低于 4 mg 的剂量时肯定不够的。

\subsection{贴片又有什么问题？}

\begin{itemize}
  \item 相比来说很贵。通常比口服药物贵。
  \item DIY 的话很难获取，一般只能通过灰色市场购买。
  \item 通常需要抗雄（参见第\ref{AA}章《抗雄激素》）。
  \item 可能引起皮肤刺激。
  \item 必须24小时贴着。
  \item 容易掉。
  \item 吸收不稳定（比如会受到温度影响）。
  \item 不容易一次性获取一大堆。
  \item 即时贴很多个也很难达到绝经前的激素水平。
\end{itemize}

\subsection{凝胶又有什么问题？}

\begin{itemize}
  \item 剂量很难精准控制，这会导致激素水平波动。
  \item 半衰期短，需要经常涂抹。
  \item 黏不拉几的。
  \item 有二手暴露的风险（接触其他人）。
  \item 通常需要抗雄（参见第\ref{AA}章《抗雄激素》）。
\end{itemize}

不过需要指出，自制凝胶需要的东西很少。对于某些人来说这是一个不错的优势。

\subsection{植入颗粒呢？}

\begin{itemize}
  \item 非常贵。
  \item 很少有人卖这个。
  \item 剂量调整周期很长。
  \item 有缺陷的颗粒会导致激素水平不足。
  \item 破碎的颗粒会导致激素水平异常升高。
  \item 基本上不适用DIY。
\end{itemize}

尤其是最后一点，这意味着你只能依赖那几个卖的很贵的药商。这可能是你第一次听说植入颗粒，这下看出问题了吧。

\subsection{那喷雾呢？}

喷雾还属于试验阶段，所以也没什么好说的。它们的优缺点和凝胶基本上相似。我提到喷雾只是想告诉你还有这种形式的雌激素存在。

\subsection{差异这么显著吗？}

\textbf{是的！}差异大到我写了这么一整篇文档，只是为了以后能少重复解释几次，为了之后能直接给对面甩个链接过去。一个剂量合理的注射方案，是我们目前能达到单药目标水平的最佳雌激素形式。

\subsection{这张表准吗？}\label{1-12}

\begin{figure}[H]
  \centering
  \includegraphics[width=1\linewidth]{STUPID_CHART_evil_bad_bad_destroy_evil_bad.png}
  \caption{狗屎表格}
  \label{fig:scebbdeb}
\end{figure}

\textbf{一点也不准。}虽然这份指南不是为了对网上每条错误信息纠错，但这张图在网上和医生提供的材料中跨语言且广泛地转播，并且造成了大量实质性伤害，因此除了在第\ref{11-6}节中简要说明外，我不得不在这里拎出来单独说说。

这张表从根本上就是错的！几乎每一项内容都在某种程度上有误导性。唯一一点正确的是：雌激素不会改变嗓音。其他的问题还有，「预计起效时间」具有误导性；「最大效果时间限制」这个概念太具有误导性了，应该把这张表的作者抓起来；大量变化（比如心理变化）没有被列出；有些内容自相矛盾，比如关于勃起功能障碍的内容；完全忽视方案的质量。

\textbf{总结来说，变化发生的速度因人而异。}这张表误导认知，歪曲事实HRT的作用，限制MtF对雌激素的理解。请彻彻底底的忘记这张表。好了，我们继续。

\section{为什么一定是注射？}

\subsection{是什么让注射这么好？}

\textbf{稳定性。} 在 HRT 中，稳定性就是一切。激素水平稳定，就意味着整体状态稳定，而稳定是好事。即便是「最差」的注射类型（后文会讲），也能比其他给药途径提供更稳定的激素周期，而这会带来很多好处。

\subsection{使用注射时还需要抗雄激素吗？}

一般来说，不需要。一个剂量合理、间隔安排得当的注射周期，只要能够持续提供足够高且稳定的雌激素水平，就可以自然地抑制睾酮分泌，从而不必使用抗雄激素——在大多数情况下这是更理想的。这被称为\textit{“单药方案（monotherapy）”}。

\subsection{单药方案是如何起作用的？}\label{2-3}

简单来讲，大脑并不在意体内具体是哪种性激素，只要够用就行。如果你体内始终有足够的激素，它就会停止进一步生产更多激素。这里的「稳定」，正是注射所能做到、而其他给药途径往往难以实现的。

更具体一点，从 \href{https://en.wikipedia.org/wiki/Hypothalamic-pituitary-gonadal_axis}{HPG轴（下丘脑—垂体—性腺轴）} 的角度来看，当血清 \textit{雌二醇}（E2）水平升高时，会抑制 \textit{促黄体生成素}（LH）与 \textit{促卵泡刺激素}（FSH）的分泌，从而抑制睾丸中的睾酮生成。

\subsection{注射为什么更安全？}

由于通常不需要抗雄激素（见第 \ref{AA} 节《抗雄激素》），因此可以避免与抗雄激素相关的长期风险。绕过肝脏代谢的生物同源雌激素（见问题 \ref{11-1}）已经尽可能接近人体自然的雌激素分泌方式，从而减少了额外风险。

\subsection{但是注射这个“物理动作”本身不是也有风险吗？}

是的，但只需进行少量训练（见第 \ref{ts} 节《注射技术与耗材》），最糟糕的情况通常也只是出现一点轻微瘀青。这有点像学骑自行车：一旦掌握了技巧，你必须非常努力才可能“严重做错”。

\subsection{注射为什么更简单？}

一旦你的方案调试完成，就可以高枕无忧。注射不需要频繁给药（例如每周一次注射 vs 每天多次服药），不像凝胶那样容易出现剂量不准确，也不会像贴片那样在周期中途脱落，更不需要像植入颗粒那样可能要长途前往医疗机构。

\subsection{注射为什么更便宜？}

简单来说，需要的雌激素总量更少。一瓶 5ml 的安瓿，通常只含约 200mg 雌激素，却足以提供接近一年的用量；而如果换算成口服药（例如 4mg × 365 天 = 1460mg），总量会大得多。这并不是严格的对比，但可以帮助理解数量级的差异。

另一个有趣的对比是：1–2 年用量的雌激素注射安瓿，体积甚至可以装进一个普通的“三个月口服药瓶”。

\subsection{但是我没有医保 / 医保不报销 / 在我这里口服更便宜 / 我所在国家没有注射剂型 / 医生不给我开注射？}

请参见第 \ref{sv} 节《药物获取》。你会大开眼界，而且很可能因此变得非常激进(?)。

\subsection{已经 HRT 好几年了，再换成注射还有意义吗？}

\textbf{有。} 虽然没有任何事情是绝对保证的，但许多人在改用注射后，即使已经 HRT 多年，仍然会体验到明显变化。例如乳房发育进一步进展、心理状态改善、抗雄激素或其他雌激素形式带来的副作用缓解、整体感觉更好等等。换用注射是值得尝试的。

\textbf{需要强调的是，在开始注射之前所经历的 HRT 时间并不是“浪费”，也不存在所谓雌激素女性化效果的“黄金窗口期”。} 注射之所以优秀，是因为它的稳定性，但最终的目标在很大程度上是相同的。参见问题 \ref{11-5} 与 \ref{11-6}。

\subsection{可是注射很可怕？}

是的，一开始确实如此。没有人喜欢针头，因为人体本能地不想被刺。但只要技术正确、耗材合适，疼痛其实非常轻微。有无数严重针头恐惧症的人，如今都觉得打针是一件无聊的小事。这种恐惧是正常且普遍的，但完全可以克服，而且值得克服。「哦，也没我想象的那么差。」是一句非常常见的话。正如那句口号所说：带着害怕去做。你会没事的。

\subsection{注射和抽血或疫苗一样吗？}

不一样。抽血通常使用更粗的针头，刺入更敏感的位置，而且还会把你的血抽走，这往往让人不舒服。疫苗之所以疼，是因为疫苗本身会引发免疫反应。而 HRT 注射只是把少量激素注入体内，而你会因为体内有激素而感觉更好。我相信你能看出区别。此外，自己给自己注射，有时甚至比让别人打更容易，这取决于个人倾向。

\subsection{有没有什么辅助注射的无障碍工具？}

有。自动注射器（auto-injector）是存在的。如果你有精细运动控制方面的问题，它可能非常有用。请参见问题 \ref{5-21}，或者继续往下读。

\subsection{但是我\textit{真的}很特殊，不能打针，比如我有玻璃骨头纸皮肤之类的\textemdash{}}

我理解这种恐惧。但如果你真的在任何情况下都不想注射，而且也没有像血友病这样的明确禁忌症，那就不要注射。直接这么说就可以。这没关系。

等你改变主意时，这份指南依然会在这里。如果你始终不想注射，那也没问题。

\section{类型与剂量}\label{td}

\subsection*{关键词汇}
\addcontentsline{toc}{subsection}{\textemdash{} 关键词汇}

\subsection{可注射雌激素有哪些不同类型？}

在 HRT 中使用的四种主要类型是 \textit{戊酸雌二醇}（EV）、\textit{环戊丙酸雌二醇}（EC）、\textit{庚酸雌二醇}（EEn）以及 \textit{十一碳酸雌二醇}（EUn）。这些都属于 \textit{雌二醇}（E2）的「酯」，在体内都会被转化为 \textit{雌二醇}。

请注意，在某些地区，口服药也会令人困惑地以 \textit{戊酸雌二醇} 的名称出售（译者注：如补佳乐/戊酸雌二醇片），但本节仅指注射剂型。

\subsection{不同类型的可注射雌激素之间有什么区别？}

各种酯之间唯一真正相关的区别在于，它们具有不同的半衰期以及由此产生的激素曲线，这反过来会影响剂量与注射频率。

\subsection{某一种可注射雌激素的女性化效果是否比另一种更好？}

\textbf{不会。} 这些差异只会影响剂量和频率，也就是体验上的定性差异，这可能会让某一种酯更符合个人偏好，但这四种类型都能很好地发挥作用，并且都保留了注射给药的优势。

\subsection{如果可以选择，我应该选哪种可注射雌激素？}

如果你有选择权，由于 \textit{庚酸雌二醇} 能提供异常稳定的激素水平，因此对大多数人来说是首选。但需要注意的是，在大多数国家，只有在 DIY 的情况下你才会拥有这个选择（见第 \ref{sv} 节《药瓶的获取》）。

如果是通过医生开具处方，你可能会有机会使用 \textit{环戊丙酸雌二醇}，但通常其浓度较低，这可能会因为需要注射较大体积而削弱其优势，这取决于你对大体积注射的耐受度。

最常被处方的可注射雌激素（尤其是在美国）是 \textit{戊酸雌二醇}。它依然完全能够带来良好效果，但存在一些小烦恼，使得在有其他选择时它并不是首选（例如 DIY 情况下）。继续往下读。

\subsection{什么是「浓度」？}

雌激素安瓿是将雌激素溶解在油性溶液中制成的。\textit{浓度}指的是该溶液中所含的雌激素量。它以质量与体积的比值表示。换句话说：每 1 毫升油（体积单位）中含有多少毫克雌激素（质量单位）。

\textbf{请理解，仅有浓度并不等于剂量。}你经常会看到用安瓿总量来表示浓度（例如 200mg / 5ml），但最好始终将其化简（本例即 40 mg/ml）。\textbf{常见浓度包括 5 mg/ml、10 mg/ml、20 mg/ml、40 mg/ml，以及偶尔出现的 50 mg/ml。}

\subsection{「剂量与频率」是什么意思？}

\textit{剂量}与\textit{频率}是决定你激素周期的两个因素。

\textit{剂量}指你一次注入体内的雌激素量（以 mg 计），而 \textit{频率}指你多久注射一次（以天或周计）。

你也会经常听到「\textit{用药方案（regimen）}」这个词，它指你正在使用的所有与 HRT 相关的药物以及各自的使用频率。

\subsection{我如何知道自己的剂量是多少？}

你的剂量等于安瓿浓度乘以你注射的体积。

\[
\text{浓度（mg/ml）} * \text{体积（ml）} = \text{剂量（mg）}
\]

\textbf{请理解，仅有体积并不等于剂量。}可以类比烘焙：你不能只说「烤 45 分钟」，因为你还需要知道烤箱设定的温度是多少。

\subsection{有哪些剂量计算示例？}

数学其实很简单，我保证！下面是一张小参考表，对比了在不同常见浓度下，为达到不同剂量所需的注射体积。只需保留两位小数即可。你使用的注射器精度并不足以支持例如 0.153ml 这样的数字，这本身就在四舍五入误差范围内，在我们的剂量尺度上没有实际意义。你也可以使用 \href{https://hrtcafe.net/Calc/}{这个计算器} 作为参考。

\begin{table}
\centering
\caption{常见浓度下按体积计算的示例剂量}
\label{tab:concentrations}
\begin{tabular}{@{}llllll@{}}
  \toprule
  \multicolumn{1}{c}{} & \multicolumn{5}{c}{浓度（mg/ml）} \\
  \cmidrule(rl){2-6}
            & 5    & 10  & 20 & 40 & 50   \\
            \cmidrule(rl){2-6}
剂量（mg） & \multicolumn{5}{c}{体积（mL）}  \\
    \cmidrule(r){1-1} \cmidrule(lr){2-6}
4        & 0.8  & 0.4 & 0.20  & 0.10 & 0.08\\
5        & 1.0    & 0.5 & 0.25 & 0.13 &  0.10\\
6        & 1.2  & 0.6   & 0.30  & 0.15 & 0.12   \\
7        & 1.4  & 0.7 & 0.35  & 0.18 & 0.14\\
8        & 1.6  & 0.8   & 0.40  & 0.20 & 0.16  \\
9        & 1.8  & 0.9 & 0.45  & 0.23 & 0.18\\
10       & 2.0    & 1.0   & 0.50  & 0.25 & 0.20 \\
  \bottomrule
\end{tabular}
\end{table}

\textbf{如何阅读这张表：} 在左侧找到你希望的剂量，然后在顶部找到你的安瓿浓度对应的列，即可得出需要注射的体积。你会注意到，对于 5 mg/ml 的安瓿，要获得合理剂量所需体积非常大。这是因为 5 mg/ml 的安瓿本身就不太理想。

\subsection{不同酯之间如何换算剂量？}

\textbf{无法换算。}由于它们的药代行为不同，并不存在一种意义上的“剂量转换”。如果你从一种酯更换为另一种，应直接从该酯的常见剂量范围开始，再根据情况调整。你可以对它们进行比较，但没有办法把一种剂量“转换”为另一种。

\subsection{如何比较不同酯的曲线与剂量？}

如果你想稍微 nerd 一点，我非常推荐 \href{http://estrannai.se}{estrannai.se}。不过这不是必须的。它是一个用于粗略比较的好工具。\href{https://estrannai.se/\#i0__cu,7,7,1-cu,5,7,3-cu,5,7,2}{这里有一个典型每周剂量的对比示例}，接下来我们会分别来看。

\textbf{需要说明的是，下面列出的剂量范围，在大多数情况下其较低端已经足够。} 从较低剂量开始，如果需要再逐步增加。剂量更多并不天然意味着更好，但这一点我们稍后会详细讨论。这些剂量范围通常不会因为你从哪里获得安瓿而发生变化。

\subsection*{认识你的酯}
\addcontentsline{toc}{subsection}{\textemdash{} 认识你的酯}

\subsection{如何给 \textit{戊酸雌二醇} 定剂量？}

为了获得良好的激素水平，使用 \textit{戊酸雌二醇} 通常需要每周一次较高剂量，或每周两次较低剂量。这取决于舒适度与耐受度。经验法则通常是：在一个周期中，每一天约对应 1mg，常见注射间隔在 3–7 天之间。\textbf{我通常推荐每周 6–8mg 的剂量}，但每 5 天 4–5mg 也非常常见。\textbf{注射间隔绝不应超过每周一次（也就是两次注射之间不应超过 7 天）。}每周一次已经接近这种酯所能维持的极限。再延长将强烈不建议，因为会增加与波动相关的副作用风险（见问题 \ref{7-3}）。

请注意，在某些地区，口服药也会令人困惑地以 \textit{戊酸雌二醇} 的名称出售，但本节仅指注射剂型。

\subsection{戊酸雌二醇的激素曲线有什么特点？}

\textit{戊酸雌二醇} 是最“难伺候”的一种酯。注射后几天内，它会迅速冲到非常高的峰值，然后同样迅速下降。这种相对不稳定性，取决于你的个体敏感性，可能会让人感觉不舒服。但通过调整频率与剂量，可以在一定程度上缓解。

\begin{figure}[H]
  \centering
  \includegraphics[width=1\linewidth]{ev.png}
  \caption{戊酸雌二醇血清雌二醇水平（pg / ml）随时间变化（天）}
  \label{fig:ev}
\end{figure}

\subsection{如何给 \textit{环戊丙酸雌二醇} 定剂量？}

\textit{环戊丙酸雌二醇} 可以轻松支持每周一次的注射。

\textbf{常见每周剂量为 5-7mg。}不建议将间隔延长至超过一周（例如每 10 天一次），因为与每周方案相比，这样使用雌激素效率更低，需要越来越高的剂量才能维持可接受水平。任何超过每周的延长都更容易因波动而产生副作用（见问题 \ref{7-3}）。

\subsection{环戊丙酸雌二醇的激素曲线有什么特点？}

\textit{环戊丙酸雌二醇} 比 \textit{戊酸雌二醇} 更缓和。其曲线变化速度较慢，高低之间的差距更小，但在典型的一周周期内仍然可以观察到明显的上升与下降。

\begin{figure}[H]
  \centering
  \includegraphics[width=1\linewidth]{ec.png}
  \caption{环戊丙酸雌二醇血清雌二醇水平（pg / ml）随时间变化（天）}
  \label{fig:ec}
\end{figure}

\subsection{如何给 \textit{庚酸雌二醇} 定剂量？}

\textit{庚酸雌二醇} 可以非常轻松地支持每周一次注射，如果愿意，甚至可以延长至约 10 天。再延长在技术上可行，但不推荐，因为激素水平会逐渐变得不稳定。\textbf{常见每周剂量为 4–6mg}，若采用最长约 10 天间隔，则推荐 5–7mg。无论如何，我仍建议每周一次，因为更稳定且更容易安排时间；将间隔延长至 10 天在我看来收益不大。

\subsection{庚酸雌二醇的激素曲线有什么特点？}

\textit{庚酸雌二醇} 是可注射雌激素中的黄金标准。其曲线在典型的一周周期内非常平缓（也就是波动很小）。这意味着可以获得极为稳定的激素水平，而不会出现与波动相关的负面副作用（见问题 \ref{7-3}）。

\begin{figure}[H]
  \centering
  \includegraphics[width=1\linewidth]{een.png}
  \caption{庚酸雌二醇血清雌二醇水平（pg / ml）随时间变化（天）}
  \label{fig:een}
\end{figure}

\subsection{如何给 \textit{十一碳酸雌二醇} 定剂量？}

\textit{十一碳酸雌二醇} 的注射间隔可以远远超过每周，甚至进入按月或按季度注射的范围。然而，目前并没有标准或公认的推荐剂量方案。影响注射后雌激素吸收的因素（即 \textit{「药代动力学」}），对于其他酯来说几乎可以忽略，但对 \textit{十一碳酸雌二醇} 却非常重要。因此，这仍然是高度实验性的领域，已经超出了本指南的范围。你可以考虑翻翻女巫的年鉴，看看满月时间表，然后在每次满月时注射一次。

\subsection{十一碳酸雌二醇的激素曲线有什么特点？}

我们其实并不清楚。现有数据过于稀少，无法完整描绘其真实曲线，而且变量非常多。如果你对此感兴趣，可以自行研究与尝试，但这属于全新的领域，你需要理解作为「人体小白鼠」的风险。因此，除非你非常清楚自己在做什么，否则我不建议这样做。

\begin{figure}[H]
  \centering
  \includegraphics[width=1\linewidth]{moon.png}
  \caption{月亮}
  \label{fig:moon}
\end{figure}

\section{血检与激素水平}

\subsection*{获取结果}
\addcontentsline{toc}{subsection}{\textemdash{} 获取结果}

\subsection{我应该多久测一次激素水平？}

在你刚开始把剂量「调到合适」的时候，你会想要相对更频繁地检测。每次调整方案之后，都应该给激素水平 1-2 个月的时间稳定下来，然后等它达到新的稳定水平后再去测一次。

\subsection{开始 HRT 之前必须先测一次激素水平吗？}

严格来说不一定，因为那时睾酮会很高、雌激素会很低，这些数据本身并不算特别有用；但出于健康考虑，仍然建议做一些常规的基础抽血检查（比如血脂全套之类）。唯一的例外是：如果你怀疑自己可能有【间性状态（intersex）】并且它可能影响你的 HRT 方案，那么开始前的血检有时能看出一些端倪。

\subsection{如果我很久没改过剂量，还需要测吗？}

严格来说也不一定，因为如果你什么都没改，那么一般也不该有什么变化。如果你更换了供货渠道/供应商、或者调整了方案的其他环节，测一次可以让你更安心；而且医生/保险也经常会要求你测。但通常不应该出现明显偏差。需要补充的是：如果你在尝试 \textit{十一碳酸雌二醇}，那无论如何你几乎都应该至少按季度（每三个月）测一次。

\subsection{我没有保险也没有医生。我去哪里做血检？}

看看你所在地区是否有自费血检/第三方检验的渠道，具体取决于当地的合规/法律环境。在很多地方，自费血检是合法的，但可能不便宜。也可能有线上渠道能让你以更低价格拿到这些检测，不过这非常依赖你所在地区的具体情况。

\subsection{我做不了/负担不起血检。我还能做 HRT 吗？}

显然，有信息总比没信息要好；但 HRT 非常安全，在常见剂量下通常不会有问题。你只是需要更依赖自己的主观感受，以及你观察到的变化。

\subsection{我应该测哪些指标？}

至少要测 \textit{雌二醇}（E2）和\textit{总睾酮}（T），因为这是最核心、最需要关注的两项。\textit{性激素结合球蛋白}（SHBG）、\textit{双氢睾酮}（DHT）、\textit{雌酮}（E1）以及\textit{催乳素}（PRL）在你遇到问题时也很值得测，因为它们有助于排查原因。\textit{促卵泡生成素}（FSH）和\textit{促黄体生成素}（LH）可以告诉你 HPG 轴是否处于被抑制/不活跃状态，这是单药治疗（monotherapy）的基础（见问题 \ref{2-3}）。但还是那句话：\textbf{\textit{雌二醇}和\textit{总睾酮}是最主要的关注点。}

\subsection{在激素周期的哪个时间点去抽血？}

在你的周期末尾（\textit{「低谷」/谷值（trough）}）。你希望尽可能接近最低点，因为这才是最有用的信息。甚至可以说：它可能是唯一有用的信息，因为「稳定的最低水平」才是主要关注点。举个例子：如果你平时是周四下午注射，那么就去下一个周四早上或中午（在下一针之前）抽血。

\subsection{医生让我测周期中点/峰值的血，应该照做吗？}

\textbf{不应该。}测雌激素峰值并不能提供有用信息，它基本只是在告诉你你用的是什么酯。往好听说，这是守旧且过时的照护标准导致的无能；往难听说，这是一种恶意：通过确保你的雌激素水平不足，来让你健康变差、见效更慢、或产生其他负面结果。\textbf{无论如何，我都建议测低谷（trough）。}

\subsection*{解读结果}
\addcontentsline{toc}{subsection}{\textemdash{} 解读结果}

\subsection{我想要的雌激素水平应该是多少？}

这大概是过渡里最有争议的问题。简短答案是：你需要足够让自己感觉良好，并且（如果你需要的话）能够抑制睾酮；除此之外，更高的水平最多只是没必要的浪费，最坏还可能适得其反。不过这个范围其实很宽，而且变量太多，个体差异永远存在。换句话说：你需要足够的雌激素让你感觉良好——就这么简单。

\subsection{雌激素越高，女性化效果会更好或更快吗？}

\textbf{不会。}高于所需的雌激素水平，可能会因为主观感受而被某些人偏好，但它不会带来额外的女性化收益。事实上，过高的水平反而可能让你感觉更糟，比如导致情绪不稳定或其他不想要的副作用。\textbf{把睾酮水平压到基线，比把雌激素水平拉到很高，对女性化更重要。}

\subsection{好吧，那我雌激素化验单上到底想看到多少？}

在理解以下前提的情况下：这个具体数字并不重要；由于检测与体内水平之间存在滞后，即使在低谷日，化验值也总会比你体内实际水平略高一点；而且这个数字会因为各种因素而落在一个「可能性云团」里——\textbf{我建议低谷（trough）至少约 200 pg/ml（730 pmol/L）。}这是一个略偏保守的建议，给你留出足够的余量，因为 HPG 轴的抑制会在低于这个水平时发生。这个区间对大多数人通常效果不错，当然也有人更偏好更高或更低。我不认为这是一个值得过度执着的数字，因为它天然会波动；如果你感觉良好，那才是最重要的。\textbf{但是}，\textbf{低谷超过 300 pg/ml（1100 pmol/L）几乎肯定已经高于所需（或应当）的水平了。}确实存在例外，但你大概率不是那个例外。总之，按你感觉最合适的来。另外也可参见问题 \ref{11-1}。

\subsection{我想要的睾酮水平应该是多少？}

抑制睾酮是获得足够女性化效果的关键要求，因此一般来说低于 50 ng/dL（1.7 nmol/L）就足够了。\textbf{值得强调的是：不追求接近 0 的睾酮。}见第 \ref{T} 章《睾酮》。

\subsection{我天生 T 很高/很低。我需要调整剂量吗？}

大概率不需要。HRT 前常见的睾酮范围几乎总是高于女性化所需的目标范围，并且无论如何都会被抑制（见问题 \ref{2-3}）。例外情况是：如果你有各种【间性相关情况】导致你需要比本指南所列建议更精细的调整——这已经超出了本指南能提供的范围。你也许不必调整，但也许你调整后感觉更好。归根结底：按你感觉最合适的来。参见问题 \ref{9-2}。

\subsection{我做过下体手术。我的雌激素水平需要不一样吗？}

因为你不再需要考虑抑制睾酮，你很可能在比现在更低的雌激素水平下也能感觉很好，\textbf{但你仍然需要雌激素。}因为你不再自己产生激素，维持足够的激素水平对你的健康至关重要。激素过低或接近没有，会带来更年期症状；这也是很多顺性别女性进入更年期后会使用 HRT 的原因。你可以按自己感觉做调整。

为了更清楚一点：\textbf{维持至少约 100 pg/ml（350 pmol/L）的最低水平，对避免骨矿物质密度（骨密度）方面的担忧是必要的。}如果你的女性化过程大体已经完成，那么在很多方面，你的激素谱会更接近一位绝经期的顺性别女性，因此她们的经验也能提供参考（见问题 \ref{11-29}）。在某些疲劳或精力不足的情况下，补充低剂量睾酮可能会有帮助（见问题 \ref{9-2}）。

\subsection{有什么因素会导致血检不准确吗？}

取决于采用的检测方法（\textit{「测定法」/assay}），生物素（biotin）补充剂可能会让雌二醇（E2）（以及其他一些指标，但我们主要关心 E2）出现异常偏高的结果。由于你不一定能提前知道会用哪种测定法，所以建议在抽血前几天暂停服用生物素补充剂。也有可能是设备或样本出了错，不过这种情况不太常见。

\subsection{不同的雌激素酯或不同给药途径，会在血检结果上表现得不一样吗？}\label{4-16}

不会。仅凭血检结果无法判断一个人用的是什么类型的雌激素。各种注射用酯都会转化成想要的雌二醇；药片、贴片、凝胶、喷雾或其他形式也是一样。说到底，都是雌激素。
\section{注射技巧与耗材}\label{ts}


\subsection*{注射部位与安全}
\addcontentsline{toc}{subsection}{\textemdash{} Sites \& Safety}

\subsection{该如何安全地进行注射？}

我推荐下面两个视频：

\begin{enumerate}
  \item \href{https://www.youtube.com/watch?v=cBabaGC2Dok}{\textit{“如何进行肌内（IM）自我注射”}}
  \item \href{https://www.youtube.com/watch?v=YfNlAZLxLyw}{\textit{“（至少对我来说目前）无痛的 IM 注射技巧”}}
\end{enumerate}

看完这两个视频，你就应该已经具备了正确注射、并把疼痛降到最低所需要的全部要点。我建议你认真学一遍，并在需要时回看。\textbf{有一件必须强调的关键细节：为了减轻疼痛，进针时要让针头斜面（bevel）朝上。}换句话说：针尖有一个非常明确的尖点，你希望最先接触皮肤的就是那个尖点。针的路径要尽可能直、干脆。如果这样更好理解，你可以想象一下手/手腕是如何旋转来完成这个动作的；但现实中，这会变成一种你自然而然学会的直觉式肌肉记忆。

\textbf{记住：注射是一项技能！}你会随着时间越来越熟练，而且并不需要很久。你可以的。

\subsection{我必须完全照这个方式注射吗？}

不必，做法有差异也完全没问题。归根结底，这件事就是「给自己扎一针」，实现它的方法有很多。找到最适合你的方式就好。通常来说，快速、像飞镖一样干脆地进针效果最好；但如果你必须慢慢来也可以，只要这种方式足够稳定，并且你能在稳定的基础上越来越熟练。

\subsection{我该如何克服注射焦虑？}

我建议把整个过程「仪式化」。当你形成一套固定流程后，它就会变成一种下意识的事情。如果你能用听音乐、聊天、看剧，或者任何能分散注意力的方式让肌肉记忆接管动作，那就更好了！找到对你有效的办法。让朋友或亲近的人帮你完成最开始的几次注射也会有帮助。对大多数人来说，第一针最可怕。通常大家最后都会说：“就这？”因为它从来都不会像他们想象得那么糟。

\subsection{我注射在身体的哪里重要吗？}

重要，也不重要。重要的是你要在安全的部位范围内注射；除此之外，你具体选哪里，主要取决于你的活动能力、单次注射的液体体积、你使用的针头/注射器组合，以及你自己的舒适度。不管怎样，\textbf{务必轮换注射部位。}每次注射都在身体左右两侧交替——例如这周打右腿，下周就打左腿。这样做是为了把长期瘢痕的风险降到最低。

\subsection{哪些注射部位是安全的？}

不同医疗机构的意见并不完全一致，而你的体型/体成分也会产生影响。我推荐像视频里那样打在大腿外侧，因为对大多数人来说都比较好操作，而且更容易做得稳定；一旦你的手法练熟了，稳定往往就意味着稳定的不疼。当然，也有人更喜欢臀部或腹部。\href{https://vertisis.com/articles/how-to-self-administer-a-subcutaneous-injection}{这个视频和网站}展示了其他一些注射部位——具体能不能用，也取决于你用的耗材。找出最适合你的方式就好。

\subsection{肌肉注射（IM）和皮下注射（SubQ/SC）是什么意思？}

在讨论注射时，你会经常听到这些词。\textit{Intramuscular} 指「肌内注射」，也就是把药液注入肌肉；\textit{subcutaneous} 指「皮下注射」，也就是把药液注入皮肤下方的脂肪层。

\subsection{肌肉注射和皮下注射什么区别？}

\textbf{在 HRT 这个语境下，肌内注射和皮下注射几乎没有区别。}皮下注射的吸收速度确实会比肌内注射更慢一些，但这种差异通常不足以影响用量。另外还要注意：一针药液很少会「完全」沉积在肌肉里或「完全」沉积在皮下层里；每一针的实际情况都会更模糊，也就进一步淡化了两者的区别。

顺带一提，药品说明书或药企资料往往会写雌二醇注射剂「仅用于肌内注射」，这通常只是因为它在审批意义上被批准的给药方式是肌注，而不是因为两者在本质上有什么不同。很多时候，不管你怎么打，小瓶里的配方其实是一样的；归根结底，重要的是瓶里的内容。继续往下看。

\subsection{我应该选择肌注还是皮下？}

\textbf{这个问题问错了。}\textbf{注射就是注射。}皮下注射之所以经常被推荐，往往是因为大家觉得「既然是皮下，就会更不疼」；但注射这件事在操作本质上并没有什么根本差别。\textbf{人们所说的那些优势，并不是注射药液沉积位置（depot）本身带来的，而是影响注射疼痛的那些因素带来的。}更好的问题应该是「我如何把注射时的疼痛降到最低？」，不过在回答它之前，还有另外两个问题要先讲。

\subsection{我的进针角度和/或偏好的注射手法重要吗？}

不重要。再强调一次：注射最重要的部分，就是让针穿过皮肤、把液体送进身体里。如果药液没有漏出来（或者至少漏得不多），而且不怎么疼（或者至少没那么疼），那你就做得非常好了。\textbf{我怎么强调都不为过：所谓「肌注 vs 皮下」的分歧根本不存在，这个问题并不会在实质上影响注射用雌激素的效果。}\textit{十一碳酸雌二醇}几乎是唯一一个「药液沉积位置」看起来会明显影响吸收的例外；但即便如此，我们也还没有完全搞清楚细节。我的意思是，请把注意力放在真正重要的事情上，而不是那些不重要的事情上。

\subsection{我需要回抽（aspirate）吗？}

不需要。Aspiration（回抽/抽吸）指的是：针扎进皮肤后、推药之前，先把活塞往后拉一下，目的是确认自己没有扎进血管。它到底有没有必要一直有争议；但对于按标准流程进行的激素注射来说，它带来的好处很少能抵过它的坏处。标准注射部位本来就不太容易碰到血管，而较短的针头长度又进一步降低了这种风险，所以大多数医疗组织已经不再推荐这样做了。

\subsection{我该如何把注射时的疼痛降到最低？}

除了练习手法、提升熟练度之外，注射不适感最主要的决定因素，是你使用的针头与注射器组合。\textbf{为了把不适降到最低，你应该在「你的药液所用载体油能推得动」的前提下，尽量选择更细的针（更高的针规），并搭配合适的注射器规格和针长。}你真正该问的是：「我应该用多细、以及多长的针来注射？」要回答这个问题，我们先讲讲针到底是怎么一回事。

\subsection*{认识你的针头}
\addcontentsline{toc}{subsection}{\textemdash{} 认识你的针头}

\subsection{针规是什么？}

\textit{针规} 用来衡量针的粗细。数字越大，针越细。例如 25G 就比 20G 更细。针规更高的针也往往更短，因为针越长就越容易弯折，所以它们的最大长度通常更小。毫不意外，更细的针一般更不疼。需要说明的是：你用的针有多粗，并不会以任何方式影响 HRT 的效果；它只会影响你注射这件事本身的难易度与舒适度。

\subsection{“Luer lock”和“insulin”注射器/针头是什么？}\label{5-13}

\textit{Luer lock（鲁尔锁）}注射器的针头和注射器本体是分开的，因此你可以分别用不同的针来「抽药」和「注射」。\textit{Insulin（胰岛素）}注射器的针是固定在注射器上的，也就是说抽药和注射都用同一根针。注意：你不想用“Luer slip（鲁尔滑口/滑接）”的针头，因为它更容易出错。在条件允许时，出于舒适度以及减少死腔（见问题 \ref{5-26}）的考虑，更推荐胰岛素注射器。

\textbf{警告：出于「扎到自己」的风险考虑，一般不建议给针回套针帽；但如果你确实要这么做（比如更换抽药针时），千万不要让你的手朝着针的方向用力。}

如果你没把针帽对准，针帽可能会裂开，从而伤到你。更推荐的做法是：把针放在水平桌面上，用「轻轻铲/套」的方式把针帽从桌面上套回去，然后把已经松松套上的针顶着墙轻压，或捏住针帽两侧把它推到位。自己给自己注射时并不存在疾病传播风险，所以你可以自行权衡是否严格遵守这条警告；但如果你给别人注射，二次扎伤（resticking）的风险就非常严重。如何处理废弃针具见问题 \ref{5-27}。

\subsection{抽药应该用多粗的针（needle gauge）？}

如果你用的是鲁尔锁注射器，建议抽药针比注射针更粗一些（也就是针规更低一些），这样从小瓶里抽药更快。但太粗也会增加落屑（刺穿胶塞时切下橡胶碎屑）的风险（见问题 \ref{5-23}），所以\textbf{建议抽药至少用 21–23G。}如果你有耐心、单次注射的体积也比较小，那么为了降低上述打芯风险，推荐用更细的针（更高的针规）。另外请注意：针在胶塞上并不会明显变钝。用胰岛素注射器时，这个问题不成立，因为针头不可拆卸。

\subsection{抽药针应该用多长？}

如果你用的是鲁尔锁注射器，抽药针的长度除了「太长会不好操作」以外，基本没太大影响。换句话说，不用纠结。用胰岛素注射器时，这个问题不成立，因为针头不可拆卸。

\subsection{注射应该用多粗的针}\label{5-16}

这是个很棘手、而且非常主观的问题。你的答案主要取决于 4 个因素：1）你注射的药液用的是什么载体油；2）瓶里是否含有共溶剂（cosolvent）；3）你能忍受针在身体里停留更久的程度；4）你愿不愿意/能不能更用力地推注射器活塞。这本质上是个舒适度问题。油越稠，用更细的针（更高针规）时就越费时、越费力；但更细的针进皮肤时也往往会明显更不疼。\textbf{作为基准，为了控制不适感，你至少应当用 25G。}大多数常见载体油一般都能比较轻松地用到 27G；其中 MCT 油尤其突出，通常能很轻松地用到 30G（见问题 \ref{6-16}）。

\subsection{注射针应该用多长？}

\textbf{I recommend between 0.5” to 1” (12.5mm to 25mm) depending on your gauge.} Below 0.5” (12.5mm) increases the likelihood of leakage. 0.25” (6.5mm) length needles can be fine depending on your technique and the fluid you’re injecting, but 0.5” (12.5mm) is a safe bet. Beyond 1” (25mm) is unnecessarily daunting and painful without any added benefits.
\textbf{我建议根据针规，选择 0.5” 到 1”（12.5mm 到 25mm）之间的针长。}短于 0.5”（12.5mm）会增加漏药的概率。0.25”（6.5mm）的针在某些手法和药液条件下也可以，但 0.5”（12.5mm）是更稳妥的选择。超过 1”（25mm）则没必要，只会更吓人、更疼，而不会带来额外好处。

\subsection{注射器的规格（容量）重要吗？}

\textbf{重要，容量大小真的很重要。}原因有二：1）容量越大的注射器刻度往往越不精细，容易导致剂量不准；2）从物理层面讲，容量越大的注射器推起来往往越费劲。为了剂量准确，你不该用一个远大于你实际注射体积的注射器（比如单次注射少于 0.1ml，就用小于 1ml 的注射器）。\textbf{能不用 3mL 注射器就尽量别用。}当然，如果你手头只有这个，那就用；但它真的不是为这种事情设计的。别问我为什么药师好像几乎只会发这个。也许是一种残酷的玩笑。

\subsection{我去哪里买注射器和针头？}

这取决于你所在地区的规定：有些地方会以禁毒为目的，禁止向个人销售针头和注射器。除此之外，医疗/兽用耗材供应商，或正规厂商的授权渠道，通常都是不错的选择。\textbf{不推荐亚马逊}，因为质量不确定。

\subsection{针头或注射器可以重复使用吗？}

\textbf{绝对不行!绝对不要重复使用任何针头或注射器，也不要和别人共用。}所有注射器——不管针头能不能拆——都只能一次性使用。你大概早就知道，但我还是要提醒你一次，因为这真的既不好也不安全。

\subsection{如果我想注射，但自己操作有困难怎么办？}\label{5-21}

你也许会喜欢自动注射器（auto-injector）。顾名思义，它能替你完成进针动作。像 \href{https://unionmedico.com/90-super-grip/}{\textit{UnionMedico 45/90 Super Grip}} 这种自动注射器可以搭配 1ml 注射器使用，能显著降低「扎进去」这一步的难度（但你仍然需要手动推活塞）；而像 \href{https://www.youtube.com/watch?v=hPbhEpUN43Y}{\textit{Owen Mumford Autoject 2}} 这种则会把胰岛素针的针头完全遮住，并自动按下活塞。网上也有各种 3D 打印的设计可用。我没有用过这些产品，这些也不构成任何背书。

\subsection*{小药瓶基础知识}
\addcontentsline{toc}{subsection}{\textemdash{} 小药瓶基础知识}

\subsection{检查小药瓶时应该看什么？}

除了留意落屑的迹象（见下文）之外，你还应该检查是否有变色、分层、污染、结晶、玻璃裂纹、纤维、头发等任何异常。合格的小药瓶不应该明显偏离正常状态。\textbf{每次使用前都要检查你的药瓶。看起来不对劲的瓶子就别用。}

\subsection{落屑是什么？}\label{5-23}

\textit{落屑（Coring）} 指的是胶塞的一小块橡胶被针头打落下来，掉进瓶里。这可能发生在：抽药针太粗、反复扎在同一个点、或者扎的次数太多（比如用很大容量的药瓶，但每次注射体积很小）。\textbf{一旦发生落屑，这个瓶子就应该立刻丢弃。}\href{https://www.youtube.com/watch?v=w5F0SLoMjC8}{\textit{45–90° 技巧}}也能帮助降低落屑的发生率。

落屑的问题在于：你不想把橡胶碎屑打进身体里。如果你能看到比较大的碎屑，那也很可能存在你肉眼看不到的更小碎屑。胶塞的作用是保护内容物不受外界影响，因此顶部被扎出洞的瓶子会更容易氧化和/或滋生细菌。\textbf{顺带一提：请确保你把新药瓶顶部的金属或塑料盖子取下来。}这听起来很显然，但有些瓶子的设计确实会让人困惑。

\subsection{药瓶多久会过期？}

未开封的药瓶如果在稳定室温、避光条件下保存，几年都不会有什么问题。假设药瓶按应有标准完成了灭菌，那么「放久了」主要担心的是载体油氧化。已经扎过的药瓶如果含有抑菌剂/防腐剂（见问题 \ref{6-17}），通常至少能用一年，或者至少能覆盖「你把它用完」所需的时间。“开封后 28 天丢弃”这类标注，只是厂商必须保证无菌性的最低期限要求，并不是最大保质期。

\subsection{我应该如何保存药瓶？}\label{5-25}

稳定室温、避光保存。高温和紫外线会导致载体油降解；低温则可能导致结晶。结晶通常可以重新溶解并混回溶液中，但如果没有完全溶解，可能会引起刺激。这一点对未开封和已开封的药瓶都一样。

另外：每次从瓶里抽出多少液体，你就应该往瓶里补回等体积的空气，避免瓶内随着时间形成真空，导致之后抽药变得非常困难。做法是：先在注射器里吸入一部分空气（体积与打算抽的药量刻度一致），把这部分空气推入药瓶，然后再按平常方式抽药。这个简单步骤能让你以后省很多麻烦。

\subsection{Dead space/死腔是什么}\label{5-26}

\textit{死腔} 指的是注射时被浪费掉的液体体积，也就是滞留在注射器或针头里的那部分药液。普通的鲁尔锁针头/注射器组合，这个数值可能高达 0.1mL；而胰岛素针的死腔则可能低到 0.003mL。出于经济性考虑，建议尽量降低死腔，因为长期累积会浪费相当多的雌激素。\href{https://hrtcafe.net/Calc/}{这个计算器}可以用来估算在不同耗材组合下会浪费多少雌激素。

还有一点要注意：如果你抽药和注射会更换针头，那么在拔掉抽药针之前，可以先把活塞稍微往回拉一点，这样抽药针里那点药液就不会被浪费掉。它的量很小，但有时也确实会有区别。\textbf{另外也要注意：注射器刻度标的是实际会被*注射进体内*的准确体积，所以你不需要为了「补偿死腔」而手动多加一点药。刻度就是准确剂量。}死腔的具体体积取决于你用的针头型号（抽药针和注射针的死腔体积通常也不会相同），所以一如既往，最重要的是保持一致性。如果你担心死腔过高，见问题 \ref{7-7} 里的另一种可能策略。

\subsection{用过的注射器和针头怎么处理？}\label{5-27}

把所有用过的注射耗材（针头、注射器等）尖端朝下放进锐器盒/利器盒（可以是专用的生物危害锐器盒，也可以重复利用一些硬塑料桶，比如蛋白粉桶或洗衣液桶）。当容器装到四分之三满时，把它封好，确保不会被意外打开。清楚标注“USED SHARPS”，然后按你所在地区的要求进行处置。\textbf{注意：锐器绝对不要丢进普通垃圾桶或回收桶。}你所在的城市/州/地区很可能会有一个网页说明家庭有害垃圾应该如何、以及在哪里处置。美国的话可以看这里：\href{https://safeneedledisposal.org/}{you can go here.}

\section{药物获取}\label{sv}

\subsection{我要去哪里弄到注射用雌激素？}

大体而言，你有两个选择：\textit{正规渠道}和 \textit{DIY 渠道}。\textit{正规渠道}通常需要医生的处方，因为在大多数国家，雌激素不能在柜台直接买到（即便买得到，一般也不包含药瓶）。\textit{DIY 渠道}则涵盖其余的一切。

\subsection{我应该用正规渠道还是 DIY 渠道？}

选择权在你，但有时你根本无从选择。两者各有利弊。当然，没人阻止你同时从多个来源获取雌激素，以兼得两者之长。在很多情况下，这甚至是值得推荐的做法。

\subsection*{正规渠道}
\addcontentsline{toc}{subsection}{\textemdash{} 正规渠道}

\subsection{正规渠道有哪些优点？}

\begin{itemize}
  \item 通常可以信任其质量控制流程与认证；
  \item 医保可能部分或全额报销；
  \item 看你遇到大夫的运气，有时会更加方便；
  \item 产品大概率批次一致稳定；
  \item \textbf{如果你在申请手术的保险批准，则可能必须至少“看起来”是在使用药物来源。}
\end{itemize}

\subsection{正规渠道有哪些缺点？}

\begin{itemize}
    \item 可供选择的酯类较少，甚至完全没有选择余地；
    \item 可能需要等待很长时间（数月甚至数年）；
    \item 可能需要医生开具处方（视所在国家而定）；
    \item 医保可能无法报销全部费用，甚至完全不予报销；
    \item 在你所在的国家，医生可能根本不会为你开具处方；
    \item 医生可能会毫无充分理由地拒绝为你开处方；
    \item 医生可能会毫无充分理由地拒绝为你续开处方；
    \item 药品短缺可能导致处方根本无法配药；
    \item 你很可能需要满足 WPATH 的严格要求，甚至更加严苛的标准；
    \item 不好囤积药物；
    \item 药物获取受到所在国家政治局势的影响。这也意味着，你的跨性别身份很可能会被记录在医疗档案中。
\end{itemize}

\subsection*{DIY 渠道}
\addcontentsline{toc}{subsection}{\textemdash{} DIY 渠道}

\subsection{DIY 渠道有哪些优点？}

\begin{itemize}
    \item 在大多数地区，通常便宜得多；
    \item 在世界各地均有获取途径；
    \item 与医疗系统的排队等待相比，获取药物所需的时间可能短数月甚至数年（唯一需要等待的通常是生产和运输）；
    \item 容易囤积；
    \item 可以选择全部酯类；
    \item 无需与医疗系统打交道；
    \item 这些药物很可能是其他人用爱制作出来的。
\end{itemize}

\subsection{DIY 渠道有哪些缺点？}

\begin{itemize}
  \item 几乎可以肯定，生产环境不符合认证洁净室的标准；
    \item 药品质量可能因来源不同而参差不齐；
    \item 根据来源不同，获取过程可能不太方便；
    \item 需要对药品来源相当信任；
    \item 需要自行寻找可靠的药贩子；
    \item 药贩子比当地药房更有可能停止运营；
    \item 产品的交付时间可能存在较大差异；
    \item 很可能需要使用加密货币付款，非常麻烦；
    \item 如果你原本可以使用医保，那么通过 DIY 渠道购买就用不了医保。
\end{itemize}
另外，如前文所述，如果你希望获得保险对手术费用的批准，保险公司很可能要求你已经持有 HRT 处方并接受治疗达到一定时长。这一点对你而言是否构成问题，取决于你的具体情况。

\subsection{哪些注射用雌激素只能通过 DIY 渠道获取？}

主要是\textit{庚酸雌二醇（estradiol enanthate）}和\textit{十一酸雌二醇（estradiol undecylate）}。正规渠道的医生几乎总是给你开 \textit{戊酸雌二醇}，但其浓度不一定能达到 40 mg/ml。医生偶尔也会开\textit{环戊丙酸雌二醇}，但浓度很少能超过 5 mg/ml 或 10 mg/ml，这会给剂量调整带来些许不便。仅仅是\textit{庚酸雌二醇}所带来的诸多优势，就已经足以成为考虑 DIY 渠道的充分理由。不过，通过 DIY 渠道，你可以获取浓度为 40 mg/ml 的任何一种酯类。对于那些想要自行尝试不同方案的人来说，\textit{十一酸雌二醇}同样只能通过 DIY 渠道获取。不过，正如前文所述，我不建议缺乏经验的人尝试使用它。

\subsection{DIY 渠道实际上都包含哪些来源？}

DIY 渠道包括商业制剂生产者、互助项目、你的朋友，以及如果你具有创业精神的话——你自己！

\subsection{我在哪里能弄到 DIY 药物？}

警哥别钓了，我才不会说呢。反正这也不是本指南的重点。已经有其他资源提供了相关信息。请专注于本指南本身。

\subsection{为什么 DIY 渠道能做到比正规渠道更便宜？}

生产一瓶药物的成本大约为 10 美元，其中包括人工成本以及按比例分摊的资本性支出（CAPEX）。这个成本还很可能高估了。商业 DIY 供应方的主要成本来自为了保持匿名性而产生的多层运营开销和运输费用。非商业性质的 DIY 来源通常不存在这些额外开销。而正规渠道通常没有任何降低价格的动力。

\subsection{DIY 合法吗？}\label{6-11}

在包括美国在内的大多数地区，雌激素并不属于受管制药物；而睾酮则可能因所在地不同而受到刑事法律的限制，也可能不受限制。在睾酮的管制问题上，美国是一个比较特殊的例外。其他国家通常不会将持有睾酮本身定为犯罪，不过，考虑到类固醇的广泛流通，即使存在相关法律，实际遭到起诉的情况也很少见。

\textbf{本指南不构成法律建议。}

\subsection{DIY 安全吗？}

作为一个广义的药物来源类别，DIY 本身既不能简单地被归类为安全，也不能简单地被归类为不安全。不过，并非所有 DIY 来源都具有相同的可靠性。当我们讨论将某种物质安全地注射进自己体内时，真正需要考虑的问题是：你是否信任生产这瓶药物的人，确信他们正确地遵循了无菌技术和操作规范，从而确保药瓶中含有你所需要的物质，而不含任何其他不该存在的东西？

对于正规药品渠道而言，在法律法规确实得到遵守的前提下，这种信任通常是与生俱来的。而对于 DIY 渠道，这种信任必须通过以下方式逐步建立：展示或解释生产流程；由独立第三方进行浓度和纯度检测；在相关社群中积累良好的声誉。

\subsection{如何判断一个 DIY 药品来源是否值得信任？}

相信你的直觉，也运用你的理性判断。

\begin{itemize}
    \item 对方是否愿意与你讨论其生产流程，或者是否在某处公开了相关信息？（例如，是否会过滤灰尘？答案必须是肯定的！！！）
    \item 对方是否表现出足够的专业能力？
    \item 他们的产品是否经过检测？
    \item 对方是否是社群中值得信赖的成员？
    \item 你信任的其他社群成员是否对其进行过审核或背书？（例如，实地检查、评价、用户反馈等）
    \item 人难免会犯错，但对方是否愿意承担责任，还是试图压制负面评价？
    \item 如果是商业供应方，他们是否会妥善处理客户订单出现的问题？
    \item 如果是商业供应方，他们是否会在产品尚未生产完成时就收取货款，却不明确告知这是预订或延期交付？（绝对不要接受这种预订！）
    \item 药瓶中是否含有防腐剂？
    \item 他们已经从事生产多长时间了？（出于充分的理由，他们可能不会透露这一信息！）
    \item 他们的产量有多大？（同样，出于充分的理由，他们可能不会透露这一信息！）
    \item 你是否总觉得哪里不太对劲？
\end{itemize}

以上只是帮助你判断一个来源是否值得信任的诸多问题中的一部分。归根结底，你需要确认的是：对方是否像你一样重视其产品的质量。

\subsection{是否应该对不同类型的 DIY 来源采用不同的审核标准？}

很可能应该如此。如果商业制剂生产者向你收取费用并提供产品，那么就应该对他们提出较高的标准，因为他们有能力将生产流程做好。另一方面，如果某个互助项目免费分发药物，那么你可能并没有太多挑选的余地。不过，这并不意味着它所提供的产品就一定更好或更差。至于你的朋友，或者你自己，那就只有你自己能够做出判断了！

\subsection*{药瓶的构造}
\addcontentsline{toc}{subsection}{\textemdash{} 药瓶的构造}

\subsection{拿到一瓶药物时，应该关注哪些方面？}

药瓶中的成分可以分为\textit{「活性成分」（active）}和\textit{「辅料」（excipient）}两类。在我们的语境中，活性成分就是雌激素酯，而辅料则是除此之外的所有其他成分。通常情况下，一瓶药物包含三种或四种成分：

\begin{enumerate}
    \item 雌激素酯；
    \item 载体油；
    \item 防腐剂；
    \item 可选的助溶剂（cosolvent）。
\end{enumerate}

我们已经在第 \ref{td} 节「类型与剂量」（TYPES AND DOSAGES）中介绍过雌激素酯。正规药品渠道提供的药瓶几乎总是包含上述全部四种成分。

\subsection{应该选择什么样的载体油？}\label{6-16}

这取决于个人偏好、个体耐受性，以及可能存在的过敏情况。\textbf{对你而言，最重要的变量是黏度，因为它会直接影响注射时的舒适度和便利性。}正如前文所讨论的（参见问题 \ref{5-16}），黏度较低的油液通常更适合搭配较大针规（更细）的针头使用，在抽取和注射时也不容易遇到困难。\textbf{HRT 中最常用的载体油是蓖麻油和 MCT 油（中链甘油三酯油）。}蓖麻油是常用载体油中黏度最高的一种，但它通常也最不容易引起刺激，因此正规药品生产商通常会使用蓖麻油。MCT 油则是常用载体油中黏度最低的一种，但有些人发现它比其他油类更容易引起刺激，而且目前主要只能通过 DIY 渠道获取。棉籽油和葡萄籽油偶尔也会被使用，但通常不是 HRT 药品生产商的首选。其他油类，例如葵花籽油、芝麻油等，也偶尔会被使用，但通常不被推荐。根据你的具体情况，这个问题可能并不重要，你可能根本没有选择余地，也可能对载体油有非常严格的要求。

\subsection{应该选择什么样的防腐剂？}\label{6-17}

注射用药瓶中最常见的防腐剂是低浓度的\textit{苯甲醇（benzyl alcohol，BA）}。这一点是强制性的，没有讨论余地：\textbf{绝对不要使用不含防腐剂的药瓶。}对于少数对苯甲醇过敏的人来说，\textit{氯丁醇（chlorobutanol）}是另一种常见的替代防腐剂。不过，DIY 渠道几乎不会使用这种防腐剂，因此可能需要专门寻找特定配方的正规药物。

\subsection{应该选择什么样的助溶剂？}

最常见的助溶剂是\textit{苯甲酸苄酯（benzyl benzoate，BB）}，它能够降低最终溶液的黏度。从技术上来说，助溶剂并非必需成分，但通常建议使用，以保证不同批次产品的一致性。在许多情况下，根据所使用的载体油和目标药物浓度，助溶剂甚至是必不可少的。有些人会觉得它具有刺激性，但也有人不会出现这种反应。

\section{TROUBLESHOOTING}

\subsection*{Dosage Uncertainty}
\addcontentsline{toc}{subsection}{\textemdash{} Dosage Uncertainty}

\subsection{My levels aren’t what I expected them to be. Why not?}

There are a number of possibilities. Recall first that model estimations cannot take into account any plethora of factors which may cause some deviation. Recall as well that it takes multiple injections until you reach stability, so if you just changed your dosage that may be why. Quadruple check with a friend that you are injecting as much as you think you are. That is more commonly an issue than you might think, but for DIY sources it is also possible that the concentration is lower than advertised due to inexperience or less precise equipment. In that case, injecting you may just need to inject a little more for that vial. \textbf{But remember, the most important thing is how you feel, not your levels.} Please note that even professional compounding pharmacies can produce dud vials not caught by quality control, as hopefully rare as that may be!

\subsection{Can I compare levels across different tests if I didn’t test at trough?}

\textbf{No.} Not accurately, anyway. This is part of why you should always test at trough. Hours before your normal time for your next injection; that’s what you want. Eliminating as many variables as possible makes the data far more useful to you. If there is nothing else that you take from this guide, please just test at trough.

\subsection{I feel really bad on my trough days. What should I do?}\label{7-3}

In most cases, either the dosage is too low or the frequency is too low. This is most pressing for \textit{estradiol valerate} and \textit{estradiol cypionate}. Adjust your dosage within the range listed or adjust the frequency. Find what works for you. It is also possible with \textit{estradiol valerate} in particular that your dosage might actually be too *high* instead of too low as the high level variability across your cycle may be the culprit for this crashing sensation. In short: swap to \textit{estradiol enanthate} if you can.

\subsection*{Injection Woes}
\addcontentsline{toc}{subsection}{\textemdash{} Injection Woes}

\subsection{The injection is harder to do when it’s cold. What should I do?}

Warm up the vial before drawing, then warm the syringe before injecting. Rolling the barrel of the syringe between your hands should be plenty to warm up the fluid. Forming this as a habit all the time should improve your injection consistency.

\subsection{The injection hurts more when it’s cold. What should I do?}

Warm up your injection site before injecting. Relaxing the muscles with a massage or a hot shower (specifically: increasing the temperature with the water aimed at your injection site before you get out) before injecting can help.

\subsection{I bled after my injection. Will I die?}

No. This means that you likely just hit a vein or a capillary which can happen sometimes. You might experience some light bruising or increased soreness. Using a cute bandage will make it heal faster.

\subsection{There was some air in my syringe. Will I die?}\label{7-7}

No. While you obviously do not want to inject just air and it can affect dosage if there is too much in the syringe, a small amount of air under 0.1ml is almost certainly not going to cause issue for you. It might actually be recommended in some cases. For instance, the \textit{air lock technique} (a standard technique for injecting fluids that are irritating or can stain; not crucial knowledge for HRT) generally involves injecting 0.1-0.3ml of air, so you have nothing to be worried about. You aren’t doing intravenous injections.

\subsection{Some of the fluid leaked out. Was my injection wasted and/or will I die?}

No. Leakage can happen for any number of reasons and is rarely enough to make a difference, so you do not need to do another injection. For the future, make sure to leave the needle in for 5-10 seconds before retracting and then apply pressure afterwards. You might consider using the air lock technique mentioned above or \href{https://www.nurse.com/nursing-resources/definitions/what-is-z-track-method/}{the Z-track method} if you are particularly concerned about leakage.

\subsection{Sometimes I am really sore after an injection. Will I die?}

No. Assuming you have otherwise followed all of the suggestions within this guide, sometimes the deposit of fluid hits an uncomfortable place for one reason or another. Better luck next time. \textbf{Make sure you alternate injection spots!} You do not want scar tissue to build up over the long term, and if a spot is already sore, you do not want to make it more sore.

\subsection{I am experiencing a lot of itchiness and irritation after injecting. Will I die?}

Probably not. There are a number of possible causes. Infection is the most concerning cause, but is unlikely in most cases. \textbf{Immediately go to a doctor if you are experiencing a fever, severe pain, muscle aches, pus, red streaks, or other signs of infection.} In most cases however, irritation like itchiness, redness, light swelling, warmth, etc are the result of using a vial whose estrogen and oil have separated (“crashed out of solution”). See below. It’s possible that you may be having a reaction to the carrier oil, but if you are suddenly experiencing issues after some injections without any issue, it is most likely that the vial contents are out of solution.

\subsection{My vial has crystals in it. Can I still use it?}

It most likely means your vial got too cold. Warm it up and gently shake to reincorporate. If the crystals are not going away, then it’s possible the vial contents have separated entirely. With a lot more heat and stirring the crystals might reincorporate, but it is simplest and safest to replace the vial if you can.



\section{PROGESTERONE}

\subsection{Do I want to take progesterone?}

\textbf{Probably.} This is a controversial question for some reason. Detractors (namely, doctors) will argue that there’s no studies to show that it plays a role in feminization therefore it should not be taken. Aside from transfeminine subjects being woefully understudied, heuristically speaking, progesterone is a key female sex hormone that plays an important role in the brain and in many functions throughout the body. Regardless of physical feminization, it is an important hormone for good health that should not be lightly overlooked.

\subsection{What is the difference between “progesterone” vs “progestin” / ”progestogen”?}

The class of hormones, both natural and synthetic, that activate the progesterone receptor are “proges\textbf{togens}”. The natural, bioidentical, and most important progestogen is “proges\textbf{terone}”. Synthetic progestogens are “proges\textbf{tins}”. These three terms are mistakenly used interchangeably in scientific literature and in clinical settings, likely causing much of the broader confusion regarding the role of progesterone in HRT, despite the fact that they are \textbf{not} equivalent.

\subsection{Do I want progesterone or a progestin?}

Progesterone. You want bioidentical progesterone.

\subsection{What’s wrong with progestins?}

Progestins, most typically \textit{medroxyprogesterone}, \textit{medroxyprogesterone acetate}, or \textit{levonorgestrel}, are generally associated with the negative side effects and long term risks (breast cancer, blood clots, depression, etc) that are falsely attributed to progesterone. They are not bioidentical which means they do not behave the same as progesterone and thus cannot be directly compared.

\subsection{What does progesterone do for feminization?}

Progesterone is believed to play a role in breast development and libido in particular, but as mentioned it’s a key hormone aside from its outward appearance effects. It does also have some antigonadotropic (i.e., it contributes to testosterone suppression) properties which can be sometimes relevant.

\subsection{Does it matter when I start progesterone?}

It is unknown. There is some belief that starting too early may harm breast development long term, but this is purely theoretical and contrary anecdotal evidence makes the answer unclear. The conservative estimate is waiting roughly a year into HRT (until Tanner Stage 3 or 4) in the possible chance that it does matter.

\subsection{How is progesterone normally taken?}

Aside from topical applications, the main form is via a pill. It is prescribed as an oral pill but is most effective when taken as a suppository. Topical sprays and creams can also work very well.

\subsection{Are you serious that progesterone should be taken as a suppository?}

Progesterone metabolizes entirely differently when taken orally vs rectally due to passing through the liver when taken orally. Oral progesterone primarily converts to \textit{allopregnanolone} which can cause heavy drowsiness, whereas rectal progesterone primarily converts to progesterone itself which is what we want (although some still converts). Some people take additional oral progesterone as a sleep aid, but please note that too much \textit{allopregnanolone} can sometimes lead to negative mental health side effects.

\subsection{How do I take progesterone as a suppository?}

Just a bit of water on the pill should work, then dry off and wash your hands. Obviously, don’t go to the bathroom for the next hour or so, so doing it before bed is best. If you are having issues with it not dissolving then you can try piercing the capsule but usually should be no issue. Be aware that if you use large homebrew suppositories made using coconut oil, the large volume of coconut oil will not want to stay in you.

\subsection{How much progesterone should I take?}

For pills, Standard dosage is 100-200mg daily at night. It is a rather arbitrary dosage; 200mg is the max that most doctors will prescribe. Some people take more than 200mg on occasion, but be aware that spiking your levels may lead to an unpleasant crash. See question below.

For topical applications, nobody can tell you with certainty due to the high variability of the delivery medium, nor is there any clear guidance on desired levels, or even frequency (likely daily), as progesterone is simply understudied. Because of this, I would advise titrating your dosage so that you understand how progesterone affects you.

\subsection{Is there any benefit to “cycling” progesterone?}\label{8-11}

No. Some people do this to mimic a cis woman’s menstrual cycle, but there is no reason to believe there is any benefit to this and it may cause negative PMS symptoms. The only exception is if you have good reason to suspect that you have an intersex condition involving a uterus that you are managing. I discourage it otherwise. See Question \ref{11-10}.

\subsection{How long should I take progesterone for?}

For as long as you plan to take estrogen and for as long as you want to. So, probably forever.

Sometimes people (or doctors) arbitrarily say to only take progesterone for X years. There is zero theoretical or empirical reason to suggest that this is sound advice. It's about as coherent as if someone (or a doctor) asked how long a trans person planned to take HRT for\textemdash{}oh wait never mind they do ask that.

\subsection{Can progesterone convert into \textit{dihydrotestosterone} (DHT)?}

No. Well, strictly speaking yes, but also no. It is largely a myth, although \href{https://whsah.co/posts/rethinking-progesterone-and-androgens/}{as outlined in detail by alix [sic] in this article}, for cases of people with \textit{nonclassical congenital adrenal hyperplasia} (ncCAH) progesterone can cause some negative side effects of increased androgenic activity. In those cases, discontinuing progesterone is recommended along with seeking out a formal diagnosis/treatment for potential adrenal disorders.

\subsection{Is there any benefit to topical progesterone applications in addition to pills?}

Maybe. It's a possible alternative to pills, especially in the case of someone with a peanut allergy since the most common pill manufacturer uses peanut oil, but again dosage is unclear. Some people find more progesterone fun, if nothing else. Be safe and have fun.

For clarity: Apply creams to your inner thigh region (elsewhere if directed), or optionally on scrotal skin (it's thin and highly vascular) in the case of sprays. And no, applying progesterone to your breasts directly is unlikely to make them grow bigger or faster compared to otherwise.

\subsection{Can I snort progesterone powder?}

Please don’t. It’s hell on your sinuses. It isn’t hard to make your own topical progesterone spray and there are guides out there. Do that instead. It’s significantly more effective, consistent, and safer.

\subsection{Where can I get progesterone?}

Progesterone tends to be more expensive through DIY sources due to the higher mass of hormones required, so ideally get it through pharmaceutical sources covered by insurance. There is also the option of grey market foreign pharmacies, which are simply pharmacies in another country, although these often require some hurdles to purchase from. Topical progesterone creams are available OTC in some locations, although it is not always the most economical depending on the concentration.

\subsection{I would like to read more about progesterone in an HRT context. What resources should I read?}\label{8-17}

Originally I linked a document here but I opted to remove it due to a number of faults that can be misleading. The problem with progesterone is that literally nobody agrees about a single aspect of it. I don't know a single source or study that people agree is good. Hell, people don't even agree if the word starts with the letter “P”. \textbf{The crucial thing to know is that progesterone is not strictly required for full feminization or good breast development, but assuming that it's not contraindicated for you, it's probably worth taking.}

It should be noted that for the entire category of progestogens there are countless myths and falsehoods invented whole cloth by both proponents and detractors alike which does not make discerning truth from the already-sloppy scholarship any easier. Fantastical claims of magical benefits and fearmongering of alleged risks based on nothing are both equally unhelpful, although the latter is worse in my opinion when comes from a medical authority, whether neglectful or malicious.

\subsection{Does progesterone interact with any other drugs related to HRT?}\label{8-18}

If you are taking 5$\alpha$-Reductase Inhibitors like \textit{finasteride} and \textit{dutasteride} (see Section \ref{AA} “ANTIANDROGENS”, or keep reading), these can affect how progesterone naturally breaks down into \textit{allopregnanolone} which can cause adverse mood effects in some people, irrespective of how you are taking progesterone. It is not fully clear how much the administration route for the 5$\alpha$-Reductase Inhibitors (i.e., topical vs oral) makes a difference, but lower systemic absorption via topical application may mitigate these side effects. It is recommended to not take either of those if you are someone affected by this interaction, but it is not in all cases anyway. Note that these depressive effects may be felt for up to a month after stopping.


\section{TESTOSTERONE}\label{T}

\subsection{Why don’t we want zero testosterone?}

Testosterone is an essential sex hormone which plays a key role in your health and well-being. We want to suppress it for feminization, but near-zero testosterone (less than 10 ng/dl, or 0.35 nmol/L) can cause issues such as poor libido, low energy, low strength (fatigue beyond just the strength loss of HRT), poor concentration, trouble sleeping, etc. Notably, issues very similar to having too little estrogen. Cis women also have more than zero testosterone, so that need not be the fear. \textbf{Adequate hormone levels are important!}

\subsection{Are there ever cases where I would want to supplement testosterone?}\label{9-2}

Yes. If you are experiencing the issues of the above and your estrogen levels are otherwise good, it’s possible that you might want to supplement with a microdose of testosterone. If you wanted to improve your erectile function, minimize any atrophy before bottom surgery, or otherwise wanted to experiment with your hormones to see what feels best for you, then that might be a reason to explore testosterone in a different context that you can hopefully appreciate more compared to pre-HRT.

\subsection{If I wanted to supplement testosterone, how would I do it?}

There’s a few possibilities. Testosterone comes in either injections or topical gels/creams, similar to estrogen as already discussed. Topical is more likely what you are going to be prescribed. Topical applications have the downsides that we have discussed for estrogen, but those are less of a concern here when precise levels are less important.

\subsection{What are the topical forms of testosterone?}

There is gel and cream. Gel is typically what will be prescribed, but some compounding pharmacies are able to make low-penetrating cream if someone wanted just topical application on the genitals. The latter is harder to get and generally more expensive, however.

\subsection{Does it matter where I apply the testosterone?}

It depends on if you have gel or cream. If you have the kind of localized cream as mentioned above, you would apply it as directly as mentioned. Otherwise, shoulders and upper arms are where gel should go. Make sure not to touch things until long after it dries!

\subsection{How much and how often should I apply testosterone?}

Season to taste. This largely depends on how you are feeling. If you have too much, you might start to experience side effects of testosterone (e.g., oily skin and body hair), but only you can say what is preferred for you. A weekly injection of 5-10mg of \textit{testosterone cypionate} might work for you, but in the case of 1\% topical gels which are often disbursed in 25/50mg packets, there is more variability. You almost never want even half a packet, and definitely not daily. I would suggest starting with much less than you think to see how you feel.

\subsection{Where would I get testosterone?}

If you are an American, you would have to get a prescription or ask any juicer at your closest Planet Fitness. Elsewhere, it depends on what gym chain is closest to you. Disclaimer: This is a joke. See Question \ref{6-11}.

\subsection{Are other steroids equivalent to testosterone in an HRT context?}

Anabolic-androgenic steroids, i.e., drugs that are structurally similar to testosterone, are not all equivalent. Commonly used black market steroids like \textit{trenbolone acetate} have a laundry list of undesirable side effects, but steroids like \textit{nandrolone decanoate} are occasionally used for postmenopausal cis women due to their relatively low androgenic properties which make them very favorable for transfeminine individuals. Regardless, in America it is unlikely you will be prescribed anything other than testosterone itself, if you are able to get a prescription at all.

\subsection{What is the relationship between testosterone and \textit{dihydrotestosterone} (DHT)?}

\textit{Dihydrotestosterone} is primarily synthesized from testosterone via the 5$\alpha$-Reductase enzyme with around 5\% of testosterone in your body being converted. Generally speaking, if testosterone levels are suppressed (or if you have had bottom surgery) then there should not be much left to convert, but systemic levels won’t be zero because it is still locally produced. Depending on your body, this would be the main reason that you might want to consider supplementing with a 5$\alpha$-Reductase Inhibitor antiandrogen as discussed in the following section. As a reminder, \textit{dihydrotestosterone} is the hormone that is responsible for body hair and hair loss.

\textbf{For any trans mascs reading this,} I will make a brief detour to note that at time of writing it is not clear what role the hormone plays with bottom growth regarding speed or total size as it relates to 5$\alpha$-Reductase inhibition. That is to say: it is known that \textit{dihydrotestosterone} plays a primary role in penile development, but it’s not clear how directly the lack thereof affects a trans masc person. Applying knowledge of micropenis treatment, we know that a topical cream is more effective than exogenous injections particularly with how \textit{dihydrotestosterone} cream is useful when a patient doesn’t respond to testosterone (particularly in the case of 5$\alpha$-Reductase deficiencies). So, food for thought. Someone get Oliver Longdick to handle the rest of this.



\section{ANTIANDROGENS}\label{AA}

\subsection{What are “antiandrogens”?}

\textit{Antiandrogens}, commonly also referred to as “T blockers” or just “blockers”, as the name(s) may suggest prevent androgens (that’s what testosterone is) from acting on your body. There are many types of antiandrogens and they are commonly prescribed as part of an HRT regimen. They are needed if someone still produces testosterone and is not doing a form of HRT conducive to monotherapy, such as injections, but they are usually not desirable. It also should be noted that (most) antiandrogens do not reduce testosterone levels in any way that matters but instead simply reduce/negate effects on the body. This is relevant when interpreting lab results and such.

\subsection{Why wouldn’t I want antiandrogens?}

The main issue with most antiandrogens is that they generally have very undesirable side effects that are superfluous if testosterone is suppressed in the first place by having enough estrogen, so those side effects are being experienced despite\textemdash{}in most cases, at least\textemdash{}being rendered unnecessary by a reasonably-dosed monotherapy regimen. Bottom surgery of any kind also makes antiandrogens unnecessary in most cases.

\subsection{When might I want antiandrogens?}

If you are not most cases, if you desire peace of mind, or if your insurance requires a prescription on file before they will cover a procedure, then you may want antiandrogens. The medications used as antiandrogens might have other effects that may be desirable outside of their antiandrogen properties depending on your health situation. Additionally, if you are supplementing androgens, you may want a \textit{dihydrotestosterone} (DHT) blocker to minimize side effects related to body hair and hair loss, but be aware that this may not be the case if you are not using bioidentical testosterone (e.g. \textit{nandrolone decanoate}) because not all androgens behave the same.

\textbf{It should be noted that temporarily using antiandrogens at the start of HRT when planning to perform monotherapy is unlikely to be necesssary nor is it recommended.} There is an adjustment period that you will experience regardless while your body adapts to the change in your primary hormones, so there is no need to overcomplicate what you are doing. Don't worry about it.

\subsection{What kinds of antiandrogens are there?}

The main medications taken as general testosterone blockers in an HRT context are \textit{spironolactone}, \textit{bicalutamide}, and \textit{cyproterone acetate}. The main medications taken to block the conversion of testosterone into \textit{dihydrotestosterone} (DHT) called “5$\alpha$-Reductase Inhibitors” (5-ARI) are \textit{finasteride} and \textit{dutasteride}. There are also GnRH agonists like \textit{leuprolide} and \textit{triptorelin}, but both of those are more often used as puberty blockers in minors, although in parts of Europe they are used for adults as well.

\subsection{When might I want to take \textit{spironolactone}?}

Due to the heroic dosages and significant negative side effects required for it to function as an antiandrogen in most cases, the only time I would ever recommend taking \textit{spironolactone} would be if you would benefit from its other effects such as its antimineralocorticoid (i.e., blocking \textit{aldosterone}) properties as it relates to blood pressure management or edema. \textbf{If you insist on taking \textit{spironolactone}, please do not take more than 100mg daily.} It has a bad reputation for a reason. “The Devil”, as it were.

In case you are unfamiliar, some of the many side effects include: brain fog, lethargy, poor memory, increased urination frequency, low blood pressure, low sodium / electrolyte imbalance, etc. In other words, \textit{spironolactone} is a blood pressure lowering diuretic that is a mediocre antiandrogen which is typically prescribed at high dosages in an otherwise-healthy population for questionably-effective off-label use. In any other healthcare context this would (or SHOULD!) be highly unadvisable given the undesirable side effect profile and the widely-available preferable alternatives that already exist, but that's the state of trans healthcare for you.

\subsection{When might I want to take \textit{bicalutamide}?}

If you are going to take an antiandrogen, \textit{bicalutamide} is likely the one to take. It is generally well tolerated, barring 1\% cases of abnormal liver function test results and symptoms of liver dysfunction, but otherwise performs the job with relatively minimal side effects. \textbf{If you take \textit{bicalutamide}, ensure regular liver function tests to make sure that your results are in range.} The liver risks are dependent on your body rather than cumulative so any problem would likely present itself within the first year. Otherwise, there should be no issues.

\subsection{When might I want to take \textit{cyproterone acetate}?}

Likely never. Take \textit{bicalutamide} instead.

The long term risk profile is poor and there is no situation that I can think of in which I would recommend this over an alternative solution. You can do everything \textit{cyproterone acetate} can by just taking more estrogen and adding progesterone to your regimen.

\subsection{When might I want to take \textit{dutasteride}?}

If you are extremely concerned about possible hair loss and/or want to maximize your chances for hair regrowth, you may want to take \textit{dutasteride}. If your testosterone is otherwise suppressed then it theoretically shouldn’t have much benefit as your \textit{dihydrotestosterone} levels should be relatively low, but bodies can be complicated, so it may be something of interest to you. Also, see Question \ref{11-14}.

It should be noted that \textit{dutasteride} can cause adverse mood effects in some people, in which case stopping is strongly recommended. Note as well that these depressive effects may be felt for up to a month after stopping.

\subsection{When might I want to take \textit{finasteride}?}

If \textit{dutasteride} is not something prescribed to you or if your insurance mandates \textit{finasteride} specifically to cover a hair treatment. Otherwise, \textit{dutasteride} is preferred as it is more effective and better tolerated.

It should be noted that \textit{finasteride} can cause adverse mood effects in some people, in which case stopping is strongly recommended. Note as well that these depressive effects may be felt for up to a month after stopping.

\subsection{Where can I get antiandrogens?}

Aside from being prescribed them by your doctor or perhaps available over-the-counter, there is also the option of grey market foreign pharmacies. These are simply pharmacies in another country, although these often take some hurdles to purchase from. \textit{Dutasteride} and \textit{finasteride} are generally the easiest to get over-the-counter because of their commonality as hair loss medication.



\section{MYTHS AND MISCS}\label{MM}

\subsection*{Common Questions}
\addcontentsline{toc}{subsection}{\textemdash{} Common Questions}

\subsection{Should I be worried about blood clots?}\label{11-1}

Yes and no. It is true that there is a correlation between estrogen dosages/levels and blood clot risk, but this is primarily related to the route of administration and the type of estrogen. Synthetic estrogens are the rightful cause of scorn and do lead to significantly increased blood clot risk, but bioidentical estrogens are not as concerning. In particular, the route of administration makes a major difference. Oral bioidentical estrogen passes through the liver which is what causes the increased blood clot risk. Injections bypass the liver, and there's no evidence to suggest nor reason to believe that injections of bioidentical estrogen provide any significant risk increase beyond the innate differences between testosterone and estrogen. The pervasive fearmongering towards all estrogen has persisted for decades despite these differences.

\textbf{If you are undergoing surgery, please know that pausing hormones out of concern for blood clots is no longer recommended by WPATH.} Many surgeons still include it in their pre-surgery guidelines out of concern for blood clots, but this is torture that has been disproven and even WPATH doesn't recommend it anymore. Remarkable, I know. Per \href{https://www.tandfonline.com/doi/pdf/10.1080/26895269.2022.2100644}{WPATH SOC 8 Statement 12.19}: \blockquote{After careful examination, investigators have found no perioperative increase in the rate of VTE [KT: \textit{venous thromboembolism}, i.e. a blood clot] among transgender individuals undergoing surgery, while being maintained on sex steroid treatment throughout when compared with that among patients whose sex steroid treatment was discontinued preoperatively (Gaither et al., 2018; Hembree et al., 2009; Kozato et al., 2021; Prince \& Safer, 2020).} I should put this in another question entirely, but to not break links, it would have to be at the bottom of a section and I think this is too important for that, so I note it here. A very important clarification that I should have had sooner.

\subsection{Is it okay to use nicotine while on HRT?}\label{11-2}

This is related to the above question. \textbf{Nicotine usage on HRT, especially if you’re on pills, compounds your risk of a blood clot on top of all the other reasons that nicotine is not good.} This extends to all forms of nicotine usage, but obviously smoking is by far the worst. You really do not want a blood clot. Even if you are not on pills, nicotine disrupts the way estrogen is metabolized and can lead to significantly reduced feminization effects. This aspect is understudied but community anecdotal reports are common. It’s not easy to quit, but I believe in you. There are good resources out there and strategies like tapering down by using patches really do work. You got this.

However, to be abundantly clear, \textbf{this does not mean that you cannot or should not take estrogen. The downsides of not taking estrogen at all far exceed the downsides of using nicotine.} This section is simply seeking to make you aware of any increased risks and potentially slower transition as a very strong recommendation and encouragement to quit. One step at a time.

\subsection{Is there benefit to starting at a low dosage vs a high dosage?}

To the best of knowledge, no. Sex hormones are not like other drugs that need to be titrated to manage side effects as we know the dosages that work for the majority of people, so personally I view “starter dosages” and “antiandrogen first” regimens as medical abuse. Some people believe that mimicking the slow timeline of puberty might be best (even though there are far more things happening than just estrogen levels), but there’s no evidence to support this. An orchiectomy day one might be best for all we know, but who is going to do that the moment they decide they are trans and/or want to start HRT?

Reframing this in another way: \textbf{there is no reason to believe that “starting slowly” on a dosage below the typical range is advantageous or preferable for feminization outcomes.} There isn't a concern of going “too fast” or anything like that. Both doctors and other trans women seemingly invent new myths by the day.

\subsection{Does body weight affect dosage?}

No. Because there is no “optimal” blood level for estrogen and because the therapeutic range of acceptable levels is so wide, \textbf{body weight does not meaningfully affect dosage for HRT.} It is for the same reason that slight deviations in dosage are unlikely to affect how you feel. There is no such thing as being “too light” or “too heavy” for HRT in any capacity.

\textbf{For the same reason, adjusting your dosage in increments of 0.1mg (note that says milligrams, not milliliters) is a difference that should not be expected to be perceived simply because our bodies are not sensitive enough to such exact measurements,} let alone the high possibility of imprecision when performing an injection that makes that certainty of this measurement unlikely. In other words, the accuracy of your dosage is more important than the precision.

\subsection{Is there such a thing as starting estrogen too late?}\label{11-5}

\textbf{No.} No matter when you start, estrogen is able to do a LOT and a proper regimen will be able to have powerful results. Sex hormones are some of the strongest hormones in our body in terms of our appearance. Everybody always wishes that they could’ve started sooner, but that’s no reason not to start now. Even if you’ve been on estrogen for years, there is still a benefit to be had in improving the quality of your regimen.

\subsection{Does feminization / breast development stop after X years?}\label{11-6}

\textbf{No.} There is not an arbitrary time where estrogen suddenly stops working. Various numbers are given and usually it’s either 1) entirely made up or 2) pointing to a study that only went for X years. Doctors in particular love to tell trans women not to expect more than B cup breasts (which isn’t even how breast sizing \textit{works}, but I digress) or for any growth after 2 years, but this is simply not true. There are cases of people who restarted estrogen after stopping for many years and still experienced new growth.

\subsection{I haven’t seen any changes in years on injections. Would swapping back to pills make a difference?}

Maybe, but maybe not. There are some anecdotes of people swapping back from injections to pills (or adding pills on top of injections) and experiencing more breast growth after “stalling out”, but the mechanism is not clear. There is speculation that the E1:E2 ratio (\textit{estrone} : \textit{estradiol}) heavily weighted towards E1 with oral pills compared to E2 for injections might make a difference for some people, although \textit{estrone} is not typically associated with feminization. There likely are other factors at play, but you are free to experiment if you wish. Data is limited.

\subsection{Is low energy and low libido normal on HRT?}

Generally, no. How libido is expressed changes in the beginning, but the vast majority of the time that someone experiences abnormally low libido it’s because they haven’t gotten their hormones sorted. The same goes for low energy. Get your hormones squared away, and barring that, check your diet/vitamins next. Make sure you don’t randomly have critically low vitamin D levels or something like that. It happens more often than you think.

\subsection{I hear about [random drug / strategy] that my friend said helps feminization. Does it actually?}

Maybe, but probably not. There is a lot of wild speculation about ways to achieve feminization goals, but many of them are akin to snake oil or have potentially serious risks far beyond HRT itself. You have the right to bodily autonomy and I cannot stop you, but I can encourage you to be smart about what you are doing. The more you get into the weeds of biology as it relates to transition, the shakier the ground becomes as quality data becomes less and less available. Desperation can lead to a lot of unwise and dangerous decisions. So be smart, and be safe.

\subsection{Do we want to mimic the estrogen cycle of cis women?}\label{11-10}

Arguably no. This is controversial, but I am of the belief that because we (well, most of us) do not have a uterus and corresponding menstrual cycle synced to our hormone levels, then there is no reason we should strive to copy that behavior. This is an \textit{is-ought} problem, in my view. The primary hormonal concern for most trans women is testosterone suppression which necessitates consistently high enough levels (barring post bottom surgery, where there is no testosterone to suppress), so high fluctuation and/or relatively low levels are likely to cause undue distress. You’re welcome to experiment, of course. Especially if testosterone suppression is no longer a concern for you. See Question \ref{8-11}, and see below.

\subsection{Do trans women experience periods?}

Similar to the last question, it’s important to understand what is happening. The unique hormone curve produced by your particular ester, your dosage, and your frequency can cause changes in your mood as your estrogen levels oscillate between injections. Some trans women liken this phenomenon to a period, but the underlying cause for these physiological changes is different and is usually a sign that your regimen needs tweaking so that you feel the best that you can as suffering is not virtuously feminine. Pain and discomfort are not requirements for womanhood nor should we assert ourselves based on bioessentialist arguments. The exception here are the intersex trans women who have a uterus and literally are having a period, in which case: yeah duh. See Question \ref{11-35}.

\subsection{Can too much estrogen convert to testosterone?}

\textbf{No.} Aromatase is the enzyme responsible for converting testosterone into estrogen, but there is no mechanism to convert estrogen into testosterone. This cannot happen. This is a completely false myth and you should be immediately wary of the knowledge level of anyone who says it to you. Unfortunately, it is doctors who repeat this myth the most.

\subsection{Does bottom surgery cause an increase in testosterone?}

No. This is not a thing. There is not a magic mechanism that suddenly causes testosterone to increase the moment that testicles are removed. Even if magic was stored in the balls, this simply isn’t how hormone production works. “Well, your adrenals…” They don’t work like that either. The only possible rare exception would be undiagnosed adrenal hyperandrogenism conditions that were suppressed by an antiandrogen like \textit{spironolactone} prior to surgery which might show itself after antiandrogens are ceased. Please stop repeating this myth.

\subsection{How do I prevent/revert hair loss?}\label{11-14}

Mechanically, it is pretty simple. A standard HRT regimen alone is borderline magic (don’t ask where the magic is stored) in this regard already, but the inclusion of 5$\alpha$-Reductase Inhibitors (5-ARI) as discussed in Section \ref{AA} “ANTIANDROGENS” is recommended in more extreme cases to completely halt any loss. Topical minoxodil 5\% is the only thing that works to firm up your hairline beyond hormones alone, but keep in mind that aside from miracle cases, you’re only saving dying/dormant follicles. Dead follicles don’t come back.

If this alone is insufficient for you, hair transplant technology has improved significantly. The Follicular Unit Extraction (FUE) procedure is what you want to look into. Here is where in the future I will link a guide written by an expert on getting insurance to cover that, once she writes it. This is peer pressure. Watch this space.

\subsection{Does exercise affect feminization?}

Probably. HRT causes gradual body recomposition, so you can help encourage your body to shift through exercise. Keep in mind that this process is VERY SLOW, so it is crucial that you eat enough to fuel how patient you have to be. The growth hormones from muscle stimulation via strength training also play a role in breast development, so it’s probably a good thing even aside from the rest of the obvious health benefits of exercise.

This is NOT just the writer’s barely-disguised fetish; strength training is important for your health! I mention this because a lot of trans women believe that touching a dumbbell will make them look like the hulk. I get it, but if you have no testosterone in you and you aren’t on steroids, then you aren’t going to look like that. Let alone the constant time, effort, and diligence required to even get close.

\subsection{What should I exercise then?}\label{11-16}

Cardio is useful for living which is important. Lower body exercises will fill out your hips and glutes to accentuate your figure. Upper body exercises will improve your posture and support your breasts which will make them look bigger. In other words, everything. You’re on estrogen. Have you seen cis women athletes? Exercise will feminize you.

\href{https://docs.google.com/document/d/1-NyE5EY5TTaRRMhk7HlTbKJ7HifjEsA4jlDO1qKQVl0/edit?tab=t.0}{This guide was shared with me} \textcolor{red}{(Warning: Google Docs link)} and looks to be a good starting place. I will note that there aren't particular exercises that feminize vs masculinize as bodies don't work like that, but you may wish to focus more on lower body exerices and flexibility more than the typical lifter.

\subsection{Can estrogen really cause height shrinkage?}

Yes. It is possible that it’s related to water content changes within tendons and ligaments, but it is not something that has been studied so the cause is fully speculation. Scientists: free study idea!

\subsection{Can estrogen really cause foot shrinkage?}

Yes. See above.

\subsection{Can estrogen really cause any other kinds of shrinkage?}\label{11-19}

Well, “use it or lose it” like they always say.

More seriously: without the presence of androgens such as testosterone regulating penile health via nocturnal erections and random erections, you are likely to experience far less erections without active effort. Without regular blood flow and oxidation via erections, the spongy tissue that composes the penis will atrophy over time via fibrosis. This leads to reduced firmness and possible pain during erections, potentially discouraging future erections and perpetuating the cycle of atrophy, which may cause issues during sex and/or for certain types of bottom surgery (namely: penile inversion vaginoplasty). It is not clear to what degree (if any) this is permanent, but the recommendations listed below in Question \ref{11-20} are more than sufficient for adequate mitigation of any issues. In simpler terms: it is a muscle to be exercised like any other.

\subsection*{Sexual Health}
\addcontentsline{toc}{subsection}{\textemdash{} Sexual Health}

\subsection{How do I improve erectile function on HRT?}\label{11-20}

Aside from using it regularly, ways to improve erectile function include: 1) Improving your fitness and physical health, particularly your cardiovascular ability; 2) consider medication like \textit{tadalafil} or \textit{sildenafil}; and 3) consider testosterone supplementation (see Section \ref{T} “TESTOSTERONE”). \textbf{Erectile dysfunction is not a guarantee on HRT and is entirely mitigable if you wish.}

For context on the two medications listed: \textit{tadalafil} is taken daily which improves the ability to have an erection at any time, whereas \textit{sildenafil} is essentially an erection on demand. They each have their own dosings and possible side effects not listed here.

If you would like to read a longer explanation for how erectile function works, \href{https://stainedglasswoman.substack.com/p/how-to-maintain-your-penis-function}{this Substack article} provides a good overview of the topic.

\subsection{How do I increase cum/pre-cum volume on HRT?}

Don’t be embarrassed, it’s a common question. Sunflower lecithin and pygeum increase both of those, in addition to testosterone supplementation as mentioned above. It seems to also make a difference for vaginal wetness and arousal for those who have had bottom surgery, but data and anecdotes are limited so it’s hard to say. Otherwise just be sure you drink enough water and have your nutrition in check.

\subsection{Can I lactate on HRT?}

Yes. Domperidone, fenugreek, sunflower lecithin, ample estrogen, and ample progesterone. Get a pump. Knock yourself out.

It should be noted that domperidone has side effects and risks associated with it, and that ability to lactate does not affect breast development. Newman-Goldfarb protocols would be what you want to look into.

\subsection{Can HRT change your senses and your perceptions, i.e. smell?}

You very likely were dissociated and depressed for years prior to starting HRT. The world is more vibrant now because you are no longer dissociating 24/7. The wonders of modern medicine!

It can, however, directly change your eye prescription. That can definitely happen.

\subsection{Can HRT change your sexuality?}

Similar to being dissociated as with above, HRT often incurs a lot more openness and acceptance with yourself which can cause a shift in how your sexuality presents itself. It is largely a semantics argument as to whether that is chemical or behavioral. A matter of perspective.

\subsection{Should I be on PrEP?}\label{11-25}

\textbf{Yes.} \href{https://en.wikipedia.org/wiki/Pre-exposure_prophylaxis_for_HIV_prevention}{\textit{Pre-exposure prophylaxis for HIV prevention} (PrEP)} is a category of antiviral drugs for the purpose of preventing HIV/AIDS. This is not directly related to HRT, but it is common for trans women to be at elevated risk of HIV/AIDs. Given the history of the AIDS pandemic, PrEP is a miracle of modern medicine that should interest you. \textbf{Note: There is no effect of any PrEP drug on HRT so you are encouraged to be on PrEP.}

If you are sexually active, you should strongly consider being on PrEP. Even if you are not sexually active, trans women are at a significantly higher risk of sexual violence, so you should still strongly consider being on PrEP. In most places you are likely to be prescribed \textit{Truvada} as a once-daily pill, although if you experience nausea as a side effect you can likely swap to \textit{Descovy} with no change in effectiveness. In the US, insurance commonly does not cover \textit{Descovy} unless you have tried or say you have tried \textit{Truvada}. The novel drug \textit{lenacapavir} as a twice-yearly injection is expected to make PrEP significantly easier as access becomes more widely available, if current options are prohibitive for you.

\subsection*{Medical Malpractice}
\addcontentsline{toc}{subsection}{\textemdash{} Medical Malpractice}

\subsection{I heard that injections are actually less stable because you do them less frequently. Is that true?}

Only if you follow the dipshit WPATH SOC 8 guidelines that list a recommended regimen of \textit{estradiol valerate} or \textit{estradiol cypionate} in the range of 5-30mg every two weeks which, to be abundantly clear, you absolutely should never do in a million years. “Do no harm”, my ass.

\subsection{But my doctor said-?}

The average doctor has essentially no training in anything related to trans healthcare, and \href{https://www.endocrine.org/news-and-advocacy/news-room/2017/endocrinologists-want-training-in-transgender-care}{4/5 endocrinologists have never had any formal training in trans healthcare}. It is most likely that you are their first trans patient and that they are inexperienced in the practical elements of managing a trans patient. Even among doctors who care a lot, they are often limited by conservative standards of care that they are forced to follow which do not always align with the care best for you. See above.

Please also be aware of “trans broken arm syndrome”, aka the tendency of doctors to blame everything on HRT. If your arm is broken, it's probably not “because of those hormones”!

And I should put this as a separate question but I don't want to break the formatting: in line with medical malpractice, there is no situation in which it is reasonable for a doctor to request to see or feel your breasts to “monitor growth” or for any other reason. It is far less common these days, thankfully, but it is sexual assault and completely unacceptable.

\subsection{My doctor won’t prescribe me injections. What do I do?}

Attempt to convince them, replace them, or seek DIY sources. Do not let a gatekeeping medical establishment prevent you from receiving the appropriate care that you deserve. \textbf{The most crucial aspect of interfacing with the medical system while trans is that you have to advocate for yourself.} This is compounded with disability, ethnicity, and other afflictions that scare doctors like womanhood.

\subsection{How does HRT for menopausal cis women relate to HRT for trans women?}\label{11-29}

While we generally have different goals and crucially have very different dosage requirements, there is an immense amount of overlap in experience for trans women and menopausal cis women. Medical misogyny in the form of incompetence, dismissiveness, antagonism, and/or misinformation is something that we unfortunately both experience. It is for this reason that it is paramount to build solidarity on this front. To give an example of what I mean, \href{https://www.youtube.com/watch?v=W0XW6av2wLQ}{the first 30-40 minutes of this interview} will likely sound extremely familiar to you if you would like to raise your blood pressure. The interviewee herself notes the connection too! The WHI ruined the lives of countless women.

\subsection*{Intersexuality and Comorbidities}
\addcontentsline{toc}{subsection}{\textemdash{} Intersexuality and Comorbidities}

\subsection{What’s up with Ehlers-Danlos Syndrome?}

This connective tissue disorder doesn’t actually relate to HRT but a lot of trans people have it so congrats in case this is how you learned that you do too. Aside from general cardiovascular long term concerns to maybe look into, keep up with strength training so that your joints work. Look into that elsewhere though. See Question \ref{11-16}.

\subsection{What kind of intersex things should I keep in mind?}

Throughout this guide, I have mentioned intersex conditions vaguely. Below is a short list of things that might be useful for you to know in your travels for yourself or for a friend.

\subsection{What’s up with Klinefelter Syndrome?}

This is a relatively (considering chromosomal mutations) common intersex-related condition that some trans women might not realize that they have as the two can overlap. It generally presents as low testosterone at the start of puberty. Good for you to know the name, just in case.

\subsection{What’s up with Persistent Müllerian Duct Syndrome (PMDS)?}

Another “I’m putting this here because this might be the first time you’ve even heard of the term” intersex-related condition that can affect some trans women, however few that may be since we don’t have numbers. The possible presence of an underdeveloped uterus leads to some possible complications and oddities. You probably extra want to have progesterone to avoid uterine cancer risks.

\subsection{What's up with ovotesticular syndrome?}

This intersex condition in particular can cause early level fluctuations which made lead to confusing test results due to the presence of both ovarian and testicular tissues, either separate or combined in an \textit{ovotestis}. This presents in many different ways which HRT can interact with as you begin suppressing \textit{luteinizing hormone} (LH). A uterus may or may not be present, multiple sets of gonads could be present, and/or it could look outwardly typical.

\subsection{What’s the difference between intestinal cramps and uterine cramps?}\label{11-35}

These are commonly misattributed in early transition as a symptom of intersex conditions. Intestinal cramps are widespread and diffuse across your abdomen, whereas uterine cramps are highly concentrated in a location somewhere below your belly button and tend to be sharp stabs/contractions in rapid succession. Like the inside of your body is used as a stress ball. Very different!

\subsection{What about other intersex conditions?}

I have listed a few notable ones, but there are far more expressions and ways of testing them that go far beyond the scope of this guide. Anecdotally, prevalence is higher than average among trans people so basic familiarity with this is useful.

\subsection*{Oddball Questions}
\addcontentsline{toc}{subsection}{\textemdash{} Oddball Questions}

\subsection{Many DIY sources only take crypto. Is that required? How does that work?}

There are other guides that cover this in better depth than I can on how to use crypto safely, including some vendors who have their own guides. But yes, crypto is often required for a lot of reasons. “Crypto” means a lot of things, but using it as a currency was the original point after all. It’s mostly just a pain in the ass. Monero (XMR) is good.

\subsection{What about Selective Estrogen Receptor Modulator (SERM) drugs for nonbinary regimens?}

Some people use SERMs as a part of a transition that is not looking to feminize as much for a more androgynous look, but it’s pretty much entirely uncharted waters thus why their mention is otherwise absent from this guide. You’re on your own if that’s something you want to explore, so please be safe. I don’t personally rate them very highly as I have not seen much to suggest that they work well for how people usually think or want them to work, at least not without a lot more caveats, but obviously there are people who like them. It's just not something I feel comfortable giving recommendations for.

The various proposed nonbinary regimens are often highly individualized because they are specific to what a persons' particular goals are. All HRT should be individualized to a degree, but there is often more variation in desired outcomes when people ask about androgyny. Hormonally, it is nontrivial. Everything stated in this guide should be treated solely as a starting place if you are wanting to experiment with something more complicated, but do remember that there is much more to achieving transition goals than just hormones alone.

\subsection{Are things like “herbal HRT” or “phytoestrogens” legitimate?}

\textbf{No.} If someone is telling you they have “herbal HRT”, they are telling you they have snake oil. The only thing that is going to feminize you is estrogen, not plant estrogens. No amount of “natural” products are a replacement for estrogen itself. This isn’t a common scam and you probably already know, but just in case you run into it, now you know for sure. If it smells like bullshit, it’s probably bullshit. Unless we’re talking about bug steroids in which case yeah those are actually cool. Won’t feminize you though.

\subsection{Is the Reddit Doctor that people constantly talk about Good?}

No.

\subsection{I hear DIY estrogen is made in a bathtub. Is that true?}

No. I honestly have no idea where or why this joke started that people now take seriously, but there’s no step in any process where a bathtub would even be considered. Don’t believe everything you read online. I don’t even know what you could even theoretically do with a bathtub, unless you think estrogen vials are full of the bathwater of trans women. I don’t know why you would think that though. It’s obviously cum.

\subsection{How does HRT affect fertility?}\label{11-42}

It is important to understand that this is extremely understudied so exact figures cannot be stated, and given the seriousness of pregnancy, I urge you to practice safe sex and lean on the side of caution where possible. HRT itself can, and likely will, make you infertile eventually, but only through full suppression of \href{https://en.wikipedia.org/wiki/Hypothalamic-pituitary-gonadal_axis}{the HPG axis} (see Question \ref{2-3}) over a long time span. In other words, if you haven't had bottom surgery of any kind and you are on an HRT regimen that is less capable of HPG axis suppression (such as pills), then this is more of a consideration.

\textbf{If the HPG axis is not suppressed then it is fully possible to impregnate someone}, and the timeline for sperm maturation is long enough that this is true even after the HPG axis has been initially suppressed for \textbf{multiple months}. Please take this very seriously. Full HPG axis suppression for at minimum six months, perhaps closer to a year out of an abundance of caution, is recommended.

\subsection{Is infertility from HRT reversible?}\label{11-43}

It is theoretically possible to reverse HRT-induced infertility, assuming you weren't already infertile prior to HRT (a large assumption!), but there are not many documented cases of this so the full efficacy of fertility restoration after long-term HRT is unknown. The process would involve restarting the HPG axis with a variety of medications along with entirely stopping HRT, which would in essence require a hormonal detransition for likely six months at minimum, and even then sperm quality is not certain or guaranteed. It is not something that should be planned for, to say the least, so planning around it would be wise. A sperm bank would be recommended before or early in HRT, financially permitting, if potential biological children are a priority and if a future relationship where that is possible/desired is likely.



\section{CREATINE}

\subsection{What is creatine?}

Creatine is an organic compound in your muscles and in your brain. It recycles ADP into ATP which is important for energy production in your body, especially initial high burst applications before other energy systems take over.

\subsection{Isn’t it like a steroid or something that bodybuilders use?}

No. Bodybuilders and athletes like it because having more energy means more activity before getting tired. They aren’t the only ones who use it since it is basically the \#1 supplement in terms of things that are actually useful and are actually researched.

\subsection{How is creatine related to HRT?}

\textbf{It isn’t!} But it’s something I yell about because I think it’s good and I am tired of repeating myself because people keep asking and you’re reading this anyway, aren’t you? I love a captive audience. My standup routine is at the bottom.

\subsection{Okay well why should I take creatine then?}

What a great question! It’s good for your brain and your muscles. Creatine is often found in relatively low concentrations for many people depending on their diet, especially people who don’t eat meat. There is compelling research about various chronic fatigue and post-viral conditions (long COVID in particular) being related to depleted creatine reserves in the brain, so some people find cognitive benefits from supplementing it. It isn’t magic but it is dirt cheap so it is worth trying in my opinion.

\subsection{What are the forms?}

Just \textit{creatine monohydrate} powder is what you want. The pills tend to be low dosage and are up charging you anyway, while gummies often destroy the creatine in the creation of the gummy. A lot of brands include creatine in some sort of mix but the pure stuff is usually cheaper.

\subsection{How do I take it then?}

The general recommendation is 5-10g daily dissolved in some sort of liquid. It dissolves best in things that aren’t just water. It’s mostly flavorless, so just throw a scoop or two in your coffee or a smoothie and call it a day. It can be a little chalky or gritty depending on the quantity and the fluid.

\subsection{Does it matter when I take it?}

Not really. It doesn’t have an immediate effect like that which is why it’s silly that it’s microdosed in pre-workout mixes. Take it whenever it’s convenient for you.

\subsection{How does it work then?}

It builds up in your body to a maximum level of saturation over a week or two. Then you just maintain that and reap the rewards (of maybe feeling better).

\subsection{Do I have to do a “loading” phase of taking more at first?}

Probably not. Unless you’re on some sort of intense training time crunch or something, this probably doesn’t matter at all. Just take whatever is convenient with some regularity.

\subsection{What are the side effects?}

Slight weight gain may be possible because of increased water weight in your muscles (which to be clear is Good, so don't be alarmed). If you don’t take it with water, or if you take too much at once, you might get a tummy ache. Ouchie.

\subsection{Who shouldn’t take it?}

People with kidney issues. Not because it causes them, but because creatinine (different spelling! Creatine becomes creatinine) is used as a marker in lab tests for a number of kidney issues and supplementing might give a false positive. Just keep it in mind.

\subsection{Do you have any brand recommendations?}

No. It shouldn’t really matter. Just get whatever seems reputable and is a reasonable price. I’d give a recommendation for the one I like but when I asked the brand for affiliate link they turned me down, so their loss! No free clout.

\subsection{You seriously put creatine into this document, huh?}

Yeah it’s pretty funny. It’s not my fault that I joked about it and people told me it legitimately helped them because now I feel obligated to keep talking about it!!!



\section{CLOSING REMARKS}

If any of the following are true:

\begin{itemize}
\item you are still mad at me despite the disclaimer;

\item you spotted an issue or typo;

\item you have a clarifying question that should be put into the text;

\item you have an objection that hopefully isn’t an Uhm Ackshually;

\item you wish to sing my praises;

\item you wish to pledge fealty;

\item you wish to send tithes my way;
\end{itemize}

Then please feel free to contact me and I’ll see what we can do. Bluesky is the easiest contact point, and you can DM me for my Signal. Otherwise, thank you for reading and I hope it helps.

\textbf{If you would like to donate to support this project,} \href{https://cash.app/Katitties}{CashApp}, \href{https://ko-fi.com/katitties}{Ko-Fi}, and \href{https://account.venmo.com/u/katitties}{Venmo} all work. I appreciate it!

And lastly: \textbf{The most important thing that you can do as a trans person is to live.} For as much as this document is a manual, it is in equal measure a message to you as a trans person that your existence is a gift upon the world, your presence is a blessing on those around you, and that you deserve to be treated with respect. Even if you do nothing else, your life is a feat worth praising. Thank you.



\section*{FRIENDS OF PGHRT}\label{FOPGHRT}
\addcontentsline{toc}{section}{FRIENDS OF PGHRT}

Across this document is a scattering of links to other guides and resources. Below is a consolidation of them which will also include more links to external resources as time goes on, ideally by other trans people. For the privacy minded or noided, note that some of these are Google Docs links.

\begin{enumerate}
  \item \href{https://startwith4mgestradiolenanthateweeklyandtestatonetothreemonths.com/}{SW4EEWATAOTTM} - TL;DR for PGHRT
  \item \href{https://hrtcafe.net/}{HRT Cafe} - HRT Resource Aggregator
  \item \href{https://transfemscience.org/}{Transfeminine Science} - Informational resource for trans medical literature
  \item \href{http://estrannai.se}{Estrannai.se} - Estradiol Pharmacokinetics Playground
  \item \href{https://globoho.moe/}{Globoho.moe} - Thailand Orchiectomy Medical Tourism Travel Guide
  \item Julia's FUE Guide - COMING SOON, I'M BULLYING HER TO WRITE FASTER
  \item \href{https://docs.google.com/document/d/1-NyE5EY5TTaRRMhk7HlTbKJ7HifjEsA4jlDO1qKQVl0/edit?tab=t.0}{Sky's Feminine Figure Beginner Program} - An exercise regimen geared towards trans fems
  \item \href{https://docs.google.com/document/d/114sztSw1aVWM2pXLDl9NrHklyvewz3EmFiHiisjM71k/edit?tab=t.0}{Sky's Diet 101} - A guide towards adjusting weight in a healthy way
  \item \href{https://stainedglasswoman.substack.com/p/how-to-maintain-your-penis-function}{How to Maintain Erectile Function on HRT} - A longer form explanation on the "use it or lose it" phenomenon
  \item \href{https://docs.google.com/document/d/1DXFxzN0XTudPZez_SO61fpqncRLPH_Be_QG_8Pcz9LU/edit?pli=1&tab=t.0}{Biohax Guide Googleslop Edition} - Trans Masc DIY Guide
\end{enumerate}

\section*{ABOUT THE AUTHOR}
\addcontentsline{toc}{section}{ABOUT THE AUTHOR}

Katie Tightpussy is an award-winning author and professional trans woman with nearly a decade of experience in the field of transgender. Her accomplishments include transiferating her sex through the novel technique of cross-sex hormone injections, being physically unable to shut up, and utilizing a very fortunate set of hyperfixations as they relate to transbobulation of the humors. She spends her days in the idyllic rural countryside of Los Angeles scheming of new ways to achieve world domination and enjoys riding her bicycle. Media inquiries can reach her agent at \href{http://katietightpussy.com}{katietightpussy.com}.



\section*{DISCLOSURES}
\addcontentsline{toc}{section}{DISCLOSURES}

No robot girls were harmed in the making of this document, including any usage of generative large language models. The author does not endorse any reproduction without attribution nor scraping of this work. Leave those poor robot girls alone.

The author declares an attraction towards women and acknowledges a potential conflict of interest for the existence of more beautiful trans women in the world.



\section*{ACKNOWLEDGEMENTS}
\addcontentsline{toc}{section}{ACKNOWLEDGEMENTS}

Though the text is primarily my voice, this document would not be even half as good without the contributions, feedback, and suggestions from others involved at every step along the way. A good reminder as ever that transition is not something best done alone.

Many thanks to Q, R, RM, and S in alphabetical order for close review and generally being fun nerds to talk to; love y’all. Special thanks to CB and J for close review that also inspired some very good bits. Thanks to KG for additional intersex information. Thanks to w [sic] for additional injection resources. Thanks to BIR collectively for a plethora of crucial nerd nitpicks. Appreciation for general review from c [sic], JTP, K, S, and V. Thanks to E and S for web dev assistance. Thanks to everyone on Bluesky who encouraged me to write this up in the first place, and everyone over the years sharing knowledge. And of course: much appreciation to all HRT nerds, even when we disagree, since we’re all trying to do the best for our community where we’ve otherwise been let down. Keep up the good work everyone.

Shout out to my IB Chemistry HL teacher many years ago who quite reasonably doubted my studiousness even though I’m now putting much of that knowledge to use for the art of transsexuality; go figure.



\section*{CHANGELOG}
\addcontentsline{toc}{section}{CHANGELOG}

\noindent \href{https://github.com/Juicysteak117/pghrt/}{Source code available here on GitHub.}

\noindent Full Compilation Datetime: \DTMnow

\noindent(There aren't LaTeXML bindings for \texttt{datetime2}, \texttt{hanging}, or \texttt{hyphenat}, so the formatting is slightly ugly. If you'd really like to help me out, please write those bindings!!!)

\noindent 2025-08-20: Initial release. 15.9k words.

\noindent 2025-08-20: A lot of typos and minor verbiage tweaks. Added Question \ref{8-18}.

\noindent 2025-08-21: Typos grow on trees. Added Question \ref{5-27}.

\noindent 2025-08-21: More tweaks. Opted to remove “WHY PROG” from Question \ref{8-17}. 17.0k words.

\noindent 2025-08-22: Nitpicks, clarifications, and typos. 17.2k words.

\noindent 2025-08-24: A few more twinks sorry tweaks. 17.3k words.

\noindent 2025-08-27: How long until remaining typos are embarrasing? 17.3k words.

\noindent 2025-08-28: Reduced ambiguity in a few areas. 17.4k words.

\noindent 2025-08-29: Additional clarity for frequencies in Section \ref{td}. 17.5k words.

\noindent 2025-09-01: Sisyphus boulder meme captioned fixing typos dot png. 17.5k words.

\noindent 2025-09-07: Added donation links per request. That's very kind. 17.5k words.

\noindent 2025-09-07: Few more tweaks. Clarified an additional common progestin. 17.6k words.

\noindent 2025-09-19: Added Question \ref{4-16} plus tweaks. 17.7k words.

\noindent 2025-09-23: A wide variety of clarifications up and down the line. 18.1k words.

\noindent 2025-09-24: Added an important note about surgery to Question \ref{11-1}. 18.3k words.

\noindent 2025-09-24: “Katie my doctor told me-” It never ends. 18.5k words.

\noindent 2025-09-26: Another pass of clarification edits. Yes I should have a git diff. Sorry that I don't. I thought I'd be done by now anyway! 18.7k words.

\noindent 2025-09-30: Added some cross references for clarity. 18.7k words.

\noindent 2025-10-02: More cross references. Likely will do another pass. 18.8k words.

\noindent 2025-10-02: Added a big bold warning about recapping to Question \ref{5-13} because SOMEONE didn't watch the video smh. 18.9k words.

\noindent 2025-10-10: Added Question \ref{11-42} and Question \ref{11-43} per request. Honestly I just forgot about fertility being a thing lol. Also added the Friends Of PGHRT postword section. 19.5k words.

\noindent 2025-10-10: Added Gretchen's Version (.txt) and fixed formatting. 19.5k words.

\noindent 2025-10-11: Added permalinks to everything, yay! And finally made a git repo. Look at me being a big girl, wow. 19.5k words.

\noindent 2025-10-11: Added an external link to Question \ref{11-20}. 19.6k words.

\noindent 2025-10-14: Reduced the pithiness and expanded the usefulness of Question \ref{11-25} per repeat request because I \textit{guess} there is public health utility in doing so. 19.8k words.

\noindent 2025-10-17: Dark mode! And a font toggle off. Plus a few links. 19.8k words.

\noindent 2025-10-25: A variety of micro tweaks and the inclusion of Question \ref{1-12}. 20.2k words

\noindent 2026-01-10: Added a clarification line under Question \ref{5-25} because that should be written somewhere. 20.3k words

\noindent 2026-01-24: Slight retweak to Question \ref{5-26} to better clarify how dead space works. 20.4k words.

\noindent 2026-01-29: OKAY FINE. Elaborated substantially on Question \ref{11-19} and Question \ref{11-20}. 20.6k words.

\end{document}